\documentclass{article}

% Language setting
% Replace `english' with e.g. `spanish' to change the document language
\usepackage[english]{babel}

% Set page size and margins
% Replace `letterpaper' with `a4paper' for UK/EU standard size
\usepackage[letterpaper,top=2cm,bottom=2cm,left=3cm,right=3cm,marginparwidth=1.75cm]{geometry}

% Useful packages
\usepackage{amssymb}
\usepackage{amsbsy}
\usepackage{amsmath}
\usepackage{mathtools}
\usepackage{graphicx}
\usepackage{braket}
\usepackage[colorlinks=true, allcolors=blue]{hyperref}

\usepackage{bm}
\newcommand\Bbbbone{%
  \ifdefined\mathbbb%
    \mathbbb{1}%
  \else%
    \boldsymbol{\mathbb{1}}%
  \fi}

\title{Flow equations}
\author{Joshua Althüser}

\DeclareMathOperator{\sgn}{sgn}
\allowdisplaybreaks

\begin{document}
\maketitle

\section{Flow}
We work with the Hamiltonian 
\begin{align}
H =	&+1 \sum_{ \mathbf{q} } \sum_{ \sigma } \varepsilon( \mathbf{q} )   \hat{c}_{ \mathbf{q}, \sigma }^\dagger  \hat{c}_{ \mathbf{q}, \sigma }^{\phantom{\dagger}}  \\
	&+1 \sum_{ \mathbf{q} \mathbf{p} \mathbf{r} } U( \mathbf{q}, \mathbf{p}, \mathbf{r} )   \hat{c}_{ \mathbf{q}, \uparrow }^\dagger  \hat{c}_{ \mathbf{p}, \downarrow }^\dagger  \hat{c}_{ \mathbf{p}-\mathbf{r}, \downarrow }^{\phantom{\dagger}} \hat{c}_{ \mathbf{q}+\mathbf{r}, \uparrow }^{\phantom{\dagger}} 
\end{align}
, and the generator of the CUT
\begin{align}
\eta =+1 \sum_{ \mathbf{q} \mathbf{p} \mathbf{r} } \alpha( \mathbf{q}, \mathbf{p}, \mathbf{r} )   \hat{c}_{ \mathbf{q}, \uparrow }^\dagger  \hat{c}_{ \mathbf{p}, \downarrow }^\dagger  \hat{c}_{ \mathbf{p}-\mathbf{r}, \downarrow }^{\phantom{\dagger}} \hat{c}_{ \mathbf{q}+\mathbf{r}, \uparrow }^{\phantom{\dagger}} \end{align}
. Here, the $\ell$-dependence of the coefficients is implied. Note that, at $\ell=0$, the interaction is independent of the momenta. 

The commutator reads
\begin{align}
[H, \eta] =	&+1 \sum_{ \mathbf{q} \mathbf{p} \mathbf{r} \mathbf{s} \mathbf{t} \mathbf{u} } U( \mathbf{q}, \mathbf{p}, \mathbf{r} )  \alpha( \mathbf{s}, \mathbf{t}, \mathbf{u} )   \hat{c}_{ \mathbf{s}, \uparrow }^\dagger  \hat{c}_{ \mathbf{q}, \uparrow }^\dagger  \hat{c}_{ \mathbf{t}, \downarrow }^\dagger  \hat{c}_{ \mathbf{p}, \downarrow }^\dagger  \hat{c}_{ \mathbf{t}+\mathbf{u}, \downarrow }^{\phantom{\dagger}} \hat{c}_{ \mathbf{p}+\mathbf{r}, \downarrow }^{\phantom{\dagger}} \hat{c}_{ \mathbf{s}-\mathbf{u}, \uparrow }^{\phantom{\dagger}} \hat{c}_{ \mathbf{q}-\mathbf{r}, \uparrow }^{\phantom{\dagger}}  \\
	&-1 \sum_{ \mathbf{q} \mathbf{p} \mathbf{r} \mathbf{s} \mathbf{t} \mathbf{u} } U( \mathbf{s}, \mathbf{t}, \mathbf{u} )  \alpha( \mathbf{q}, \mathbf{p}, \mathbf{r} )   \hat{c}_{ \mathbf{s}, \uparrow }^\dagger  \hat{c}_{ \mathbf{q}, \uparrow }^\dagger  \hat{c}_{ \mathbf{t}, \downarrow }^\dagger  \hat{c}_{ \mathbf{p}, \downarrow }^\dagger  \hat{c}_{ \mathbf{t}+\mathbf{u}, \downarrow }^{\phantom{\dagger}} \hat{c}_{ \mathbf{p}+\mathbf{r}, \downarrow }^{\phantom{\dagger}} \hat{c}_{ \mathbf{s}-\mathbf{u}, \uparrow }^{\phantom{\dagger}} \hat{c}_{ \mathbf{q}-\mathbf{r}, \uparrow }^{\phantom{\dagger}}  \\
	&-1 \sum_{ \mathbf{q} \mathbf{p} \mathbf{r} \mathbf{s} \mathbf{t} } U( \mathbf{p}-\mathbf{r}, \mathbf{q}, \mathbf{p} )  \alpha( \mathbf{r}, \mathbf{s}, \mathbf{t} )   \hat{c}_{ -\mathbf{p}+\mathbf{r}, \uparrow }^\dagger  \hat{c}_{ \mathbf{s}, \downarrow }^\dagger  \hat{c}_{ \mathbf{q}, \downarrow }^\dagger  \hat{c}_{ \mathbf{s}-\mathbf{t}, \downarrow }^{\phantom{\dagger}} \hat{c}_{ -\mathbf{p}+\mathbf{q}, \downarrow }^{\phantom{\dagger}} \hat{c}_{ \mathbf{r}+\mathbf{t}, \uparrow }^{\phantom{\dagger}}  \\
	&-1 \sum_{ \mathbf{q} \mathbf{p} \mathbf{r} \mathbf{s} \mathbf{t} } U( \mathbf{q}, \mathbf{p}+\mathbf{s}, \mathbf{p} )  \alpha( \mathbf{r}, \mathbf{s}, \mathbf{t} )   \hat{c}_{ \mathbf{r}, \uparrow }^\dagger  \hat{c}_{ \mathbf{q}, \uparrow }^\dagger  \hat{c}_{ \mathbf{p}+\mathbf{s}, \downarrow }^\dagger  \hat{c}_{ \mathbf{s}-\mathbf{t}, \downarrow }^{\phantom{\dagger}} \hat{c}_{ \mathbf{r}+\mathbf{t}, \uparrow }^{\phantom{\dagger}} \hat{c}_{ \mathbf{p}+\mathbf{q}, \uparrow }^{\phantom{\dagger}}  \\
	&+1 \sum_{ \mathbf{q} \mathbf{p} \mathbf{r} \mathbf{s} \mathbf{t} } U( \mathbf{r}, \mathbf{s}, \mathbf{t} )  \alpha( \mathbf{p}-\mathbf{r}, \mathbf{q}, \mathbf{p} )   \hat{c}_{ -\mathbf{p}+\mathbf{r}, \uparrow }^\dagger  \hat{c}_{ \mathbf{s}, \downarrow }^\dagger  \hat{c}_{ \mathbf{q}, \downarrow }^\dagger  \hat{c}_{ \mathbf{s}-\mathbf{t}, \downarrow }^{\phantom{\dagger}} \hat{c}_{ -\mathbf{p}+\mathbf{q}, \downarrow }^{\phantom{\dagger}} \hat{c}_{ \mathbf{r}+\mathbf{t}, \uparrow }^{\phantom{\dagger}}  \\
	&+1 \sum_{ \mathbf{q} \mathbf{p} \mathbf{r} \mathbf{s} \mathbf{t} } U( \mathbf{r}, \mathbf{s}, \mathbf{t} )  \alpha( \mathbf{q}, \mathbf{p}+\mathbf{s}, \mathbf{p} )   \hat{c}_{ \mathbf{r}, \uparrow }^\dagger  \hat{c}_{ \mathbf{q}, \uparrow }^\dagger  \hat{c}_{ \mathbf{p}+\mathbf{s}, \downarrow }^\dagger  \hat{c}_{ \mathbf{s}-\mathbf{t}, \downarrow }^{\phantom{\dagger}} \hat{c}_{ \mathbf{r}+\mathbf{t}, \uparrow }^{\phantom{\dagger}} \hat{c}_{ \mathbf{p}+\mathbf{q}, \uparrow }^{\phantom{\dagger}}  \\
	&+1 \sum_{ \mathbf{q} \mathbf{p} \mathbf{r} \mathbf{s} } U( \mathbf{p}-\mathbf{q}, \mathbf{q}+\mathbf{r}, \mathbf{q} )  \alpha( \mathbf{p}, \mathbf{r}, \mathbf{s} )   \hat{c}_{ \mathbf{p}-\mathbf{q}, \uparrow }^\dagger  \hat{c}_{ \mathbf{q}+\mathbf{r}, \downarrow }^\dagger  \hat{c}_{ \mathbf{r}-\mathbf{s}, \downarrow }^{\phantom{\dagger}} \hat{c}_{ \mathbf{p}+\mathbf{s}, \uparrow }^{\phantom{\dagger}}  \\
	&-1 \sum_{ \mathbf{q} \mathbf{p} \mathbf{r} \mathbf{s} } U( \mathbf{p}, \mathbf{r}, \mathbf{s} )  \alpha( \mathbf{p}-\mathbf{q}, \mathbf{q}+\mathbf{r}, \mathbf{q} )   \hat{c}_{ \mathbf{p}-\mathbf{q}, \uparrow }^\dagger  \hat{c}_{ \mathbf{q}+\mathbf{r}, \downarrow }^\dagger  \hat{c}_{ \mathbf{r}-\mathbf{s}, \downarrow }^{\phantom{\dagger}} \hat{c}_{ \mathbf{p}+\mathbf{s}, \uparrow }^{\phantom{\dagger}}  \\
	&-1 \sum_{ \mathbf{q} \mathbf{p} \mathbf{r} \mathbf{s} } \sum_{ \sigma } \alpha( \mathbf{q}, \mathbf{p}, \mathbf{r} )  \varepsilon( \mathbf{s} )   \hat{c}_{ \mathbf{q}, \uparrow }^\dagger  \hat{c}_{ \mathbf{s}, \sigma }^\dagger  \hat{c}_{ \mathbf{p}, \downarrow }^\dagger  \hat{c}_{ \mathbf{p}-\mathbf{r}, \downarrow }^{\phantom{\dagger}} \hat{c}_{ \mathbf{s}, \sigma }^{\phantom{\dagger}} \hat{c}_{ \mathbf{q}+\mathbf{r}, \uparrow }^{\phantom{\dagger}}  \\
	&+1 \sum_{ \mathbf{q} \mathbf{p} \mathbf{r} \mathbf{s} } \sum_{ \sigma } \alpha( \mathbf{p}, \mathbf{r}, \mathbf{s} )  \varepsilon( \mathbf{q} )   \hat{c}_{ \mathbf{p}, \uparrow }^\dagger  \hat{c}_{ \mathbf{q}, \sigma }^\dagger  \hat{c}_{ \mathbf{r}, \downarrow }^\dagger  \hat{c}_{ \mathbf{r}-\mathbf{s}, \downarrow }^{\phantom{\dagger}} \hat{c}_{ \mathbf{q}, \sigma }^{\phantom{\dagger}} \hat{c}_{ \mathbf{p}+\mathbf{s}, \uparrow }^{\phantom{\dagger}}  \\
	&+1 \sum_{ \mathbf{q} \mathbf{p} \mathbf{r} } \alpha( \mathbf{q}, \mathbf{p}, \mathbf{r} )  \varepsilon( \mathbf{q} )   \hat{c}_{ \mathbf{q}, \uparrow }^\dagger  \hat{c}_{ \mathbf{p}, \downarrow }^\dagger  \hat{c}_{ \mathbf{p}-\mathbf{r}, \downarrow }^{\phantom{\dagger}} \hat{c}_{ \mathbf{q}+\mathbf{r}, \uparrow }^{\phantom{\dagger}}  \\
	&+1 \sum_{ \mathbf{q} \mathbf{p} \mathbf{r} } \alpha( \mathbf{q}, \mathbf{p}, \mathbf{r} )  \varepsilon( \mathbf{p} )   \hat{c}_{ \mathbf{q}, \uparrow }^\dagger  \hat{c}_{ \mathbf{p}, \downarrow }^\dagger  \hat{c}_{ \mathbf{p}-\mathbf{r}, \downarrow }^{\phantom{\dagger}} \hat{c}_{ \mathbf{q}+\mathbf{r}, \uparrow }^{\phantom{\dagger}}  \\
	&-1 \sum_{ \mathbf{q} \mathbf{p} \mathbf{r} } \alpha( \mathbf{p}-\mathbf{r}, \mathbf{q}, \mathbf{p} )  \varepsilon( \mathbf{r} )   \hat{c}_{ -\mathbf{p}+\mathbf{r}, \uparrow }^\dagger  \hat{c}_{ \mathbf{q}, \downarrow }^\dagger  \hat{c}_{ -\mathbf{p}+\mathbf{q}, \downarrow }^{\phantom{\dagger}} \hat{c}_{ \mathbf{r}, \uparrow }^{\phantom{\dagger}}  \\
	&-1 \sum_{ \mathbf{q} \mathbf{p} \mathbf{r} } \alpha( \mathbf{q}, \mathbf{p}+\mathbf{r}, \mathbf{p} )  \varepsilon( \mathbf{r} )   \hat{c}_{ \mathbf{q}, \uparrow }^\dagger  \hat{c}_{ \mathbf{p}+\mathbf{r}, \downarrow }^\dagger  \hat{c}_{ \mathbf{r}, \downarrow }^{\phantom{\dagger}} \hat{c}_{ \mathbf{p}+\mathbf{q}, \uparrow }^{\phantom{\dagger}} 
\end{align}
After normal ordering with respect to the Fermi sea, we omit any contribution with more than 4 operators. The result reads
\begin{align*}
		&+1 \sum_{ \mathbf{q} \mathbf{p} \mathbf{r} } \alpha( \mathbf{q}, \mathbf{r}, \mathbf{p} )  \varepsilon( \mathbf{r} )   \langle \hat{n}_{ \mathbf{r}, \downarrow } \rangle :\hat{c}_{ \mathbf{q}, \uparrow }^\dagger \hat{c}_{ \mathbf{r}, \downarrow }^\dagger \hat{c}_{ -\mathbf{p}+\mathbf{r}, \downarrow }^{\phantom{\dagger}}\hat{c}_{ \mathbf{p}+\mathbf{q}, \uparrow }^{\phantom{\dagger}}: \\
	&+1 \sum_{ \mathbf{q} \mathbf{p} \mathbf{r} } \alpha( \mathbf{q}, \mathbf{p}+\mathbf{r}, \mathbf{p} )  \varepsilon( \mathbf{r} )   \langle \hat{n}_{ \mathbf{r}, \downarrow } \rangle :\hat{c}_{ \mathbf{q}, \uparrow }^\dagger \hat{c}_{ \mathbf{p}+\mathbf{r}, \downarrow }^\dagger \hat{c}_{ \mathbf{r}, \downarrow }^{\phantom{\dagger}}\hat{c}_{ \mathbf{p}+\mathbf{q}, \uparrow }^{\phantom{\dagger}}: \\
	&-1 \sum_{ \mathbf{q} \mathbf{p} \mathbf{r} \mathbf{s} } \sum_{ \sigma } \alpha( \mathbf{q}, \mathbf{p}, \mathbf{r} )  \varepsilon( \mathbf{s} )   \langle \hat{n}_{ \mathbf{s}, \sigma } \rangle :\hat{c}_{ \mathbf{q}, \uparrow }^\dagger \hat{c}_{ \mathbf{p}, \downarrow }^\dagger \hat{c}_{ \mathbf{p}-\mathbf{r}, \downarrow }^{\phantom{\dagger}}\hat{c}_{ \mathbf{q}+\mathbf{r}, \uparrow }^{\phantom{\dagger}}: \\
	&+1 \sum_{ \mathbf{q} \mathbf{p} \mathbf{r} } \alpha( \mathbf{p}-\mathbf{r}, \mathbf{q}, \mathbf{p} )  \varepsilon( \mathbf{r} )   \langle \hat{n}_{ \mathbf{r}, \uparrow } \rangle :\hat{c}_{ -\mathbf{p}+\mathbf{r}, \uparrow }^\dagger \hat{c}_{ \mathbf{q}, \downarrow }^\dagger \hat{c}_{ -\mathbf{p}+\mathbf{q}, \downarrow }^{\phantom{\dagger}}\hat{c}_{ \mathbf{r}, \uparrow }^{\phantom{\dagger}}: \\
	&+1 \sum_{ \mathbf{q} \mathbf{p} \mathbf{r} } \alpha( \mathbf{r}, \mathbf{q}, \mathbf{p} )  \varepsilon( \mathbf{r} )   \langle \hat{n}_{ \mathbf{r}, \uparrow } \rangle :\hat{c}_{ \mathbf{r}, \uparrow }^\dagger \hat{c}_{ \mathbf{q}, \downarrow }^\dagger \hat{c}_{ -\mathbf{p}+\mathbf{q}, \downarrow }^{\phantom{\dagger}}\hat{c}_{ \mathbf{p}+\mathbf{r}, \uparrow }^{\phantom{\dagger}}: \\
	&-1 \sum_{ \mathbf{q} \mathbf{p} \mathbf{r} } \sum_{ \sigma } \alpha( \mathbf{q}, \mathbf{p}, 0 )  \varepsilon( \mathbf{r} )   \langle \hat{n}_{ \mathbf{q}, \uparrow } \rangle :\hat{c}_{ \mathbf{r}, \sigma }^\dagger \hat{c}_{ \mathbf{p}, \downarrow }^\dagger \hat{c}_{ \mathbf{p}, \downarrow }^{\phantom{\dagger}}\hat{c}_{ \mathbf{r}, \sigma }^{\phantom{\dagger}}: \\
	&-1 \sum_{ \mathbf{q} \mathbf{p} \mathbf{r} } \sum_{ \sigma } \alpha( \mathbf{q}, \mathbf{p}, 0 )  \varepsilon( \mathbf{r} )   \langle \hat{n}_{ \mathbf{p}, \downarrow } \rangle :\hat{c}_{ \mathbf{q}, \uparrow }^\dagger \hat{c}_{ \mathbf{r}, \sigma }^\dagger \hat{c}_{ \mathbf{r}, \sigma }^{\phantom{\dagger}}\hat{c}_{ \mathbf{q}, \uparrow }^{\phantom{\dagger}}: \\
	&+1 \sum_{ \mathbf{q} \mathbf{p} \mathbf{r} } \sum_{ \sigma } \alpha( \mathbf{p}, \mathbf{r}, 0 )  \varepsilon( \mathbf{q} )   \langle \hat{n}_{ \mathbf{r}, \downarrow } \rangle :\hat{c}_{ \mathbf{p}, \uparrow }^\dagger \hat{c}_{ \mathbf{q}, \sigma }^\dagger \hat{c}_{ \mathbf{q}, \sigma }^{\phantom{\dagger}}\hat{c}_{ \mathbf{p}, \uparrow }^{\phantom{\dagger}}: \\
	&-1 \sum_{ \mathbf{q} \mathbf{p} \mathbf{r} } \alpha( \mathbf{q}, \mathbf{p}, \mathbf{r} )  \varepsilon( \mathbf{p} )   \langle \hat{n}_{ \mathbf{p}, \downarrow } \rangle :\hat{c}_{ \mathbf{q}, \uparrow }^\dagger \hat{c}_{ \mathbf{p}, \downarrow }^\dagger \hat{c}_{ \mathbf{p}-\mathbf{r}, \downarrow }^{\phantom{\dagger}}\hat{c}_{ \mathbf{q}+\mathbf{r}, \uparrow }^{\phantom{\dagger}}: \\
	&-1 \sum_{ \mathbf{q} \mathbf{p} \mathbf{r} } \alpha( \mathbf{q}, \mathbf{p}, \mathbf{r} )  \varepsilon( \mathbf{p}+\mathbf{r} )   \langle \hat{n}_{ \mathbf{p}+\mathbf{r}, \downarrow } \rangle :\hat{c}_{ \mathbf{q}, \uparrow }^\dagger \hat{c}_{ \mathbf{p}, \downarrow }^\dagger \hat{c}_{ \mathbf{p}-\mathbf{r}, \downarrow }^{\phantom{\dagger}}\hat{c}_{ \mathbf{q}+\mathbf{r}, \uparrow }^{\phantom{\dagger}}: \\
	&+1 \sum_{ \mathbf{q} \mathbf{p} \mathbf{r} \mathbf{s} } \sum_{ \sigma } \alpha( \mathbf{p}, \mathbf{r}, \mathbf{s} )  \varepsilon( \mathbf{q} )   \langle \hat{n}_{ \mathbf{q}, \sigma } \rangle :\hat{c}_{ \mathbf{p}, \uparrow }^\dagger \hat{c}_{ \mathbf{r}, \downarrow }^\dagger \hat{c}_{ \mathbf{r}-\mathbf{s}, \downarrow }^{\phantom{\dagger}}\hat{c}_{ \mathbf{p}+\mathbf{s}, \uparrow }^{\phantom{\dagger}}: \\
	&-1 \sum_{ \mathbf{q} \mathbf{p} \mathbf{r} } \alpha( \mathbf{q}, \mathbf{p}, \mathbf{r} )  \varepsilon( \mathbf{q}+\mathbf{r} )   \langle \hat{n}_{ \mathbf{q}+\mathbf{r}, \uparrow } \rangle :\hat{c}_{ \mathbf{q}, \uparrow }^\dagger \hat{c}_{ \mathbf{p}, \downarrow }^\dagger \hat{c}_{ \mathbf{p}-\mathbf{r}, \downarrow }^{\phantom{\dagger}}\hat{c}_{ \mathbf{q}+\mathbf{r}, \uparrow }^{\phantom{\dagger}}: \\
	&-1 \sum_{ \mathbf{q} \mathbf{p} \mathbf{r} } \alpha( \mathbf{q}, \mathbf{p}, \mathbf{r} )  \varepsilon( \mathbf{q} )   \langle \hat{n}_{ \mathbf{q}, \uparrow } \rangle :\hat{c}_{ \mathbf{q}, \uparrow }^\dagger \hat{c}_{ \mathbf{p}, \downarrow }^\dagger \hat{c}_{ \mathbf{p}-\mathbf{r}, \downarrow }^{\phantom{\dagger}}\hat{c}_{ \mathbf{q}+\mathbf{r}, \uparrow }^{\phantom{\dagger}}: \\
	&+1 \sum_{ \mathbf{q} \mathbf{p} \mathbf{r} } \sum_{ \sigma } \alpha( \mathbf{p}, \mathbf{r}, 0 )  \varepsilon( \mathbf{q} )   \langle \hat{n}_{ \mathbf{p}, \uparrow } \rangle :\hat{c}_{ \mathbf{q}, \sigma }^\dagger \hat{c}_{ \mathbf{r}, \downarrow }^\dagger \hat{c}_{ \mathbf{r}, \downarrow }^{\phantom{\dagger}}\hat{c}_{ \mathbf{q}, \sigma }^{\phantom{\dagger}}: \\
	&+1 \sum_{ \mathbf{q} \mathbf{p} \mathbf{r} } \alpha( \mathbf{q}, \mathbf{p}, \mathbf{r} )  \varepsilon( \mathbf{q} )   :\hat{c}_{ \mathbf{q}, \uparrow }^\dagger \hat{c}_{ \mathbf{p}, \downarrow }^\dagger \hat{c}_{ \mathbf{p}-\mathbf{r}, \downarrow }^{\phantom{\dagger}}\hat{c}_{ \mathbf{q}+\mathbf{r}, \uparrow }^{\phantom{\dagger}}: \\
	&+1 \sum_{ \mathbf{q} \mathbf{p} \mathbf{r} } \alpha( \mathbf{q}, \mathbf{p}, \mathbf{r} )  \varepsilon( \mathbf{p} )   :\hat{c}_{ \mathbf{q}, \uparrow }^\dagger \hat{c}_{ \mathbf{p}, \downarrow }^\dagger \hat{c}_{ \mathbf{p}-\mathbf{r}, \downarrow }^{\phantom{\dagger}}\hat{c}_{ \mathbf{q}+\mathbf{r}, \uparrow }^{\phantom{\dagger}}: \\
	&-1 \sum_{ \mathbf{q} \mathbf{p} \mathbf{r} } \alpha( \mathbf{p}-\mathbf{r}, \mathbf{q}, \mathbf{p} )  \varepsilon( \mathbf{r} )   :\hat{c}_{ -\mathbf{p}+\mathbf{r}, \uparrow }^\dagger \hat{c}_{ \mathbf{q}, \downarrow }^\dagger \hat{c}_{ -\mathbf{p}+\mathbf{q}, \downarrow }^{\phantom{\dagger}}\hat{c}_{ \mathbf{r}, \uparrow }^{\phantom{\dagger}}: \\
	&-1 \sum_{ \mathbf{q} \mathbf{p} \mathbf{r} } \alpha( \mathbf{q}, \mathbf{p}+\mathbf{r}, \mathbf{p} )  \varepsilon( \mathbf{r} )   :\hat{c}_{ \mathbf{q}, \uparrow }^\dagger \hat{c}_{ \mathbf{p}+\mathbf{r}, \downarrow }^\dagger \hat{c}_{ \mathbf{r}, \downarrow }^{\phantom{\dagger}}\hat{c}_{ \mathbf{p}+\mathbf{q}, \uparrow }^{\phantom{\dagger}}: \\
	&+1 \sum_{ \mathbf{q} \mathbf{p} } \alpha( \mathbf{p}, \mathbf{q}, 0 )  \varepsilon( \mathbf{p} )   \langle \hat{n}_{ \mathbf{p}, \uparrow } \rangle \langle \hat{n}_{ \mathbf{p}, \uparrow } \rangle :\hat{c}_{ \mathbf{q}, \downarrow }^\dagger \hat{c}_{ \mathbf{q}, \downarrow }^{\phantom{\dagger}}: \\
	&+1 \sum_{ \mathbf{q} \mathbf{p} \mathbf{r} } \sum_{ \sigma } \alpha( \mathbf{p}, \mathbf{r}, 0 )  \varepsilon( \mathbf{q} )   \langle \hat{n}_{ \mathbf{r}, \downarrow } \rangle \langle \hat{n}_{ \mathbf{q}, \sigma } \rangle :\hat{c}_{ \mathbf{p}, \uparrow }^\dagger \hat{c}_{ \mathbf{p}, \uparrow }^{\phantom{\dagger}}: \\
	&-2 \sum_{ \mathbf{q} \mathbf{p} } \alpha( \mathbf{q}, \mathbf{p}, 0 )  \varepsilon( \mathbf{q} )   \langle \hat{n}_{ \mathbf{q}, \uparrow } \rangle \langle \hat{n}_{ \mathbf{p}, \downarrow } \rangle :\hat{c}_{ \mathbf{q}, \uparrow }^\dagger \hat{c}_{ \mathbf{q}, \uparrow }^{\phantom{\dagger}}: \\
	&-1 \sum_{ \mathbf{q} \mathbf{p} } \alpha( \mathbf{q}, \mathbf{p}, 0 )  \varepsilon( \mathbf{q} )   \langle \hat{n}_{ \mathbf{q}, \uparrow } \rangle \langle \hat{n}_{ \mathbf{q}, \uparrow } \rangle :\hat{c}_{ \mathbf{p}, \downarrow }^\dagger \hat{c}_{ \mathbf{p}, \downarrow }^{\phantom{\dagger}}: \\
	&+1 \sum_{ \mathbf{q} \mathbf{p} \mathbf{r} } \sum_{ \sigma } \alpha( \mathbf{p}, \mathbf{r}, 0 )  \varepsilon( \mathbf{q} )   \langle \hat{n}_{ \mathbf{r}, \downarrow } \rangle \langle \hat{n}_{ \mathbf{p}, \uparrow } \rangle :\hat{c}_{ \mathbf{q}, \sigma }^\dagger \hat{c}_{ \mathbf{q}, \sigma }^{\phantom{\dagger}}: \\
	&+1 \sum_{ \mathbf{q} \mathbf{p} \mathbf{r} } \sum_{ \sigma } \alpha( \mathbf{p}, \mathbf{r}, 0 )  \varepsilon( \mathbf{q} )   \langle \hat{n}_{ \mathbf{q}, \sigma } \rangle \langle \hat{n}_{ \mathbf{p}, \uparrow } \rangle :\hat{c}_{ \mathbf{r}, \downarrow }^\dagger \hat{c}_{ \mathbf{r}, \downarrow }^{\phantom{\dagger}}: \\
	&-1 \sum_{ \mathbf{q} \mathbf{p} \mathbf{r} } \sum_{ \sigma } \alpha( \mathbf{q}, \mathbf{p}, 0 )  \varepsilon( \mathbf{r} )   \langle \hat{n}_{ \mathbf{q}, \uparrow } \rangle \langle \hat{n}_{ \mathbf{p}, \downarrow } \rangle :\hat{c}_{ \mathbf{r}, \sigma }^\dagger \hat{c}_{ \mathbf{r}, \sigma }^{\phantom{\dagger}}: \\
	&+1 \sum_{ \mathbf{q} \mathbf{p} } \alpha( \mathbf{q}, \mathbf{p}, 0 )  \varepsilon( \mathbf{q} )   \langle \hat{n}_{ \mathbf{p}, \downarrow } \rangle :\hat{c}_{ \mathbf{q}, \uparrow }^\dagger \hat{c}_{ \mathbf{q}, \uparrow }^{\phantom{\dagger}}: \\
	&+1 \sum_{ \mathbf{q} \mathbf{p} } \alpha( \mathbf{q}, \mathbf{p}, 0 )  \varepsilon( \mathbf{q} )   \langle \hat{n}_{ \mathbf{q}, \uparrow } \rangle :\hat{c}_{ \mathbf{p}, \downarrow }^\dagger \hat{c}_{ \mathbf{p}, \downarrow }^{\phantom{\dagger}}: \\
	&-1 \sum_{ \mathbf{q} \mathbf{p} \mathbf{r} } \sum_{ \sigma } \alpha( \mathbf{q}, \mathbf{p}, 0 )  \varepsilon( \mathbf{r} )   \langle \hat{n}_{ \mathbf{r}, \sigma } \rangle \langle \hat{n}_{ \mathbf{q}, \uparrow } \rangle :\hat{c}_{ \mathbf{p}, \downarrow }^\dagger \hat{c}_{ \mathbf{p}, \downarrow }^{\phantom{\dagger}}: \\
	&-1 \sum_{ \mathbf{q} \mathbf{p} \mathbf{r} } \sum_{ \sigma } \alpha( \mathbf{q}, \mathbf{p}, 0 )  \varepsilon( \mathbf{r} )   \langle \hat{n}_{ \mathbf{r}, \sigma } \rangle \langle \hat{n}_{ \mathbf{p}, \downarrow } \rangle :\hat{c}_{ \mathbf{q}, \uparrow }^\dagger \hat{c}_{ \mathbf{q}, \uparrow }^{\phantom{\dagger}}: \\
	&-1 \sum_{ \mathbf{q} \mathbf{p} } \alpha( \mathbf{p}, \mathbf{q}, 0 )  \varepsilon( \mathbf{p} )   \langle \hat{n}_{ \mathbf{q}, \downarrow } \rangle :\hat{c}_{ \mathbf{p}, \uparrow }^\dagger \hat{c}_{ \mathbf{p}, \uparrow }^{\phantom{\dagger}}: \\
	&-1 \sum_{ \mathbf{q} \mathbf{p} } \alpha( \mathbf{p}, \mathbf{q}, 0 )  \varepsilon( \mathbf{p} )   \langle \hat{n}_{ \mathbf{p}, \uparrow } \rangle :\hat{c}_{ \mathbf{q}, \downarrow }^\dagger \hat{c}_{ \mathbf{q}, \downarrow }^{\phantom{\dagger}}: \\
	&+2 \sum_{ \mathbf{q} \mathbf{p} } \alpha( \mathbf{p}, \mathbf{q}, 0 )  \varepsilon( \mathbf{p} )   \langle \hat{n}_{ \mathbf{q}, \downarrow } \rangle \langle \hat{n}_{ \mathbf{p}, \uparrow } \rangle :\hat{c}_{ \mathbf{p}, \uparrow }^\dagger \hat{c}_{ \mathbf{p}, \uparrow }^{\phantom{\dagger}}: \\
	&+1 \sum_{ \mathbf{q} \mathbf{p} \mathbf{r} } \sum_{ \sigma } \alpha( \mathbf{p}, \mathbf{r}, 0 )  \varepsilon( \mathbf{q} )   \langle \hat{n}_{ \mathbf{r}, \downarrow } \rangle \langle \hat{n}_{ \mathbf{q}, \sigma } \rangle \langle \hat{n}_{ \mathbf{p}, \uparrow } \rangle : \hat{1} : \\
	&-1 \sum_{ \mathbf{q} \mathbf{p} \mathbf{r} } \sum_{ \sigma } \alpha( \mathbf{q}, \mathbf{p}, 0 )  \varepsilon( \mathbf{r} )   \langle \hat{n}_{ \mathbf{r}, \sigma } \rangle \langle \hat{n}_{ \mathbf{q}, \uparrow } \rangle \langle \hat{n}_{ \mathbf{p}, \downarrow } \rangle : \hat{1} : \\
	&+1 \sum_{ \mathbf{q} \mathbf{p} } \alpha( \mathbf{p}, \mathbf{q}, 0 )  \varepsilon( \mathbf{p} )   \langle \hat{n}_{ \mathbf{q}, \downarrow } \rangle \langle \hat{n}_{ \mathbf{p}, \uparrow } \rangle \langle \hat{n}_{ \mathbf{p}, \uparrow } \rangle : \hat{1} : \\
	&-1 \sum_{ \mathbf{q} \mathbf{p} } \alpha( \mathbf{q}, \mathbf{p}, 0 )  \varepsilon( \mathbf{q} )   \langle \hat{n}_{ \mathbf{q}, \uparrow } \rangle \langle \hat{n}_{ \mathbf{q}, \uparrow } \rangle \langle \hat{n}_{ \mathbf{p}, \downarrow } \rangle : \hat{1} : \\
	&-1 \sum_{ \mathbf{q} \mathbf{p} } \alpha( \mathbf{p}, \mathbf{q}, 0 )  \varepsilon( \mathbf{p} )   \langle \hat{n}_{ \mathbf{q}, \downarrow } \rangle \langle \hat{n}_{ \mathbf{p}, \uparrow } \rangle : \hat{1} : \\
	&+1 \sum_{ \mathbf{q} \mathbf{p} } \alpha( \mathbf{q}, \mathbf{p}, 0 )  \varepsilon( \mathbf{q} )   \langle \hat{n}_{ \mathbf{q}, \uparrow } \rangle \langle \hat{n}_{ \mathbf{p}, \downarrow } \rangle : \hat{1} : \\
	&-1 \sum_{ \mathbf{q} \mathbf{p} \mathbf{r} \mathbf{s} } U( \mathbf{p}-\mathbf{q}, \mathbf{q}+\mathbf{r}, \mathbf{q} )  \alpha( \mathbf{p}, \mathbf{r}, \mathbf{s} )   \langle \hat{n}_{ \mathbf{r}, \downarrow } \rangle :\hat{c}_{ \mathbf{p}-\mathbf{q}, \uparrow }^\dagger \hat{c}_{ \mathbf{q}+\mathbf{r}, \downarrow }^\dagger \hat{c}_{ \mathbf{r}-\mathbf{s}, \downarrow }^{\phantom{\dagger}}\hat{c}_{ \mathbf{p}+\mathbf{s}, \uparrow }^{\phantom{\dagger}}: \\
	&-1 \sum_{ \mathbf{q} \mathbf{p} \mathbf{r} \mathbf{s} } U( \mathbf{p}+\mathbf{s}, \mathbf{q}, \mathbf{s} )  \alpha( \mathbf{p}, \mathbf{r}, \mathbf{s} )   \langle \hat{n}_{ \mathbf{p}+\mathbf{s}, \uparrow } \rangle :\hat{c}_{ \mathbf{r}, \downarrow }^\dagger \hat{c}_{ \mathbf{q}, \downarrow }^\dagger \hat{c}_{ \mathbf{r}-\mathbf{s}, \downarrow }^{\phantom{\dagger}}\hat{c}_{ \mathbf{q}+\mathbf{s}, \downarrow }^{\phantom{\dagger}}: \\
	&-1 \sum_{ \mathbf{q} \mathbf{p} \mathbf{r} \mathbf{s} } U( \mathbf{q}, \mathbf{s}, \mathbf{p} )  \alpha( \mathbf{r}, \mathbf{s}, 0 )   \langle \hat{n}_{ \mathbf{s}, \downarrow } \rangle \langle \hat{n}_{ \mathbf{r}, \uparrow } \rangle :\hat{c}_{ \mathbf{q}, \uparrow }^\dagger \hat{c}_{ \mathbf{s}, \downarrow }^\dagger \hat{c}_{ \mathbf{p}+\mathbf{s}, \downarrow }^{\phantom{\dagger}}\hat{c}_{ -\mathbf{p}+\mathbf{q}, \uparrow }^{\phantom{\dagger}}: \\
	&-1 \sum_{ \mathbf{q} \mathbf{p} \mathbf{r} \mathbf{s} } U( \mathbf{q}, \mathbf{r}+\mathbf{s}, \mathbf{s} )  \alpha( \mathbf{p}, \mathbf{r}, \mathbf{s} )   \langle \hat{n}_{ \mathbf{r}+\mathbf{s}, \downarrow } \rangle :\hat{c}_{ \mathbf{p}, \uparrow }^\dagger \hat{c}_{ \mathbf{q}, \uparrow }^\dagger \hat{c}_{ \mathbf{p}+\mathbf{s}, \uparrow }^{\phantom{\dagger}}\hat{c}_{ \mathbf{q}-\mathbf{s}, \uparrow }^{\phantom{\dagger}}: \\
	&-1 \sum_{ \mathbf{q} \mathbf{p} \mathbf{r} \mathbf{s} } U( \mathbf{p}+\mathbf{s}, \mathbf{q}+\mathbf{r}, \mathbf{q} )  \alpha( \mathbf{p}, \mathbf{r}, \mathbf{s} )   \langle \hat{n}_{ \mathbf{p}+\mathbf{s}, \uparrow } \rangle :\hat{c}_{ \mathbf{p}, \uparrow }^\dagger \hat{c}_{ \mathbf{q}+\mathbf{r}, \downarrow }^\dagger \hat{c}_{ \mathbf{r}-\mathbf{s}, \downarrow }^{\phantom{\dagger}}\hat{c}_{ \mathbf{p}+\mathbf{q}+\mathbf{s}, \uparrow }^{\phantom{\dagger}}: \\
	&+1 \sum_{ \mathbf{q} \mathbf{p} \mathbf{r} \mathbf{s} } U( \mathbf{q}, \mathbf{r}, 0 )  \alpha( \mathbf{p}, \mathbf{r}, \mathbf{s} )   \langle \hat{n}_{ \mathbf{q}, \uparrow } \rangle :\hat{c}_{ \mathbf{p}, \uparrow }^\dagger \hat{c}_{ \mathbf{r}, \downarrow }^\dagger \hat{c}_{ \mathbf{r}-\mathbf{s}, \downarrow }^{\phantom{\dagger}}\hat{c}_{ \mathbf{p}+\mathbf{s}, \uparrow }^{\phantom{\dagger}}: \\
	&+1 \sum_{ \mathbf{q} \mathbf{p} \mathbf{r} \mathbf{s} } U( \mathbf{q}, \mathbf{p}+\mathbf{s}, \mathbf{p} )  \alpha( \mathbf{r}, \mathbf{s}, 0 )   \langle \hat{n}_{ \mathbf{r}, \uparrow } \rangle :\hat{c}_{ \mathbf{q}, \uparrow }^\dagger \hat{c}_{ \mathbf{p}+\mathbf{s}, \downarrow }^\dagger \hat{c}_{ \mathbf{s}, \downarrow }^{\phantom{\dagger}}\hat{c}_{ \mathbf{p}+\mathbf{q}, \uparrow }^{\phantom{\dagger}}: \\
	&-1 \sum_{ \mathbf{q} \mathbf{p} \mathbf{r} \mathbf{s} } U( \mathbf{p}-\mathbf{q}, \mathbf{q}+\mathbf{r}, \mathbf{q} )  \alpha( \mathbf{p}, \mathbf{r}, \mathbf{s} )   \langle \hat{n}_{ \mathbf{p}, \uparrow } \rangle :\hat{c}_{ \mathbf{p}-\mathbf{q}, \uparrow }^\dagger \hat{c}_{ \mathbf{q}+\mathbf{r}, \downarrow }^\dagger \hat{c}_{ \mathbf{r}-\mathbf{s}, \downarrow }^{\phantom{\dagger}}\hat{c}_{ \mathbf{p}+\mathbf{s}, \uparrow }^{\phantom{\dagger}}: \\
	&+1 \sum_{ \mathbf{q} \mathbf{p} \mathbf{r} \mathbf{s} } U( \mathbf{q}, \mathbf{p}, 0 )  \alpha( \mathbf{r}, \mathbf{s}, 0 )   \langle \hat{n}_{ \mathbf{r}, \uparrow } \rangle \langle \hat{n}_{ \mathbf{p}, \downarrow } \rangle :\hat{c}_{ \mathbf{q}, \uparrow }^\dagger \hat{c}_{ \mathbf{s}, \downarrow }^\dagger \hat{c}_{ \mathbf{s}, \downarrow }^{\phantom{\dagger}}\hat{c}_{ \mathbf{q}, \uparrow }^{\phantom{\dagger}}: \\
	&+1 \sum_{ \mathbf{q} \mathbf{p} \mathbf{r} \mathbf{s} } U( \mathbf{p}, \mathbf{r}, \mathbf{s} )  \alpha( \mathbf{p}+\mathbf{q}, \mathbf{r}+\mathbf{s}, \mathbf{q} )   \langle \hat{n}_{ \mathbf{r}+\mathbf{s}, \downarrow } \rangle :\hat{c}_{ \mathbf{p}-\mathbf{q}, \uparrow }^\dagger \hat{c}_{ \mathbf{r}, \downarrow }^\dagger \hat{c}_{ -\mathbf{q}+\mathbf{r}-\mathbf{s}, \downarrow }^{\phantom{\dagger}}\hat{c}_{ \mathbf{p}+\mathbf{s}, \uparrow }^{\phantom{\dagger}}: \\
	&-1 \sum_{ \mathbf{q} \mathbf{p} \mathbf{r} \mathbf{s} } U( \mathbf{p}, \mathbf{r}, \mathbf{s} )  \alpha( \mathbf{p}, \mathbf{q}, 0 )   \langle \hat{n}_{ \mathbf{q}, \downarrow } \rangle :\hat{c}_{ \mathbf{p}, \uparrow }^\dagger \hat{c}_{ \mathbf{r}, \downarrow }^\dagger \hat{c}_{ \mathbf{r}-\mathbf{s}, \downarrow }^{\phantom{\dagger}}\hat{c}_{ \mathbf{p}+\mathbf{s}, \uparrow }^{\phantom{\dagger}}: \\
	&-1 \sum_{ \mathbf{q} \mathbf{p} \mathbf{r} \mathbf{s} } U( \mathbf{r}, \mathbf{s}, 0 )  \alpha( \mathbf{p}+\mathbf{r}, \mathbf{q}, \mathbf{p} )   \langle \hat{n}_{ \mathbf{s}, \downarrow } \rangle :\hat{c}_{ -\mathbf{p}+\mathbf{r}, \uparrow }^\dagger \hat{c}_{ \mathbf{q}, \downarrow }^\dagger \hat{c}_{ -\mathbf{p}+\mathbf{q}, \downarrow }^{\phantom{\dagger}}\hat{c}_{ \mathbf{r}, \uparrow }^{\phantom{\dagger}}: \\
	&+1 \sum_{ \mathbf{q} \mathbf{p} \mathbf{r} \mathbf{s} } U( \mathbf{p}, \mathbf{r}, \mathbf{s} )  \alpha( \mathbf{p}-\mathbf{q}, \mathbf{q}+\mathbf{r}, \mathbf{q} )   \langle \hat{n}_{ \mathbf{r}, \downarrow } \rangle :\hat{c}_{ \mathbf{p}-\mathbf{q}, \uparrow }^\dagger \hat{c}_{ \mathbf{q}+\mathbf{r}, \downarrow }^\dagger \hat{c}_{ \mathbf{r}-\mathbf{s}, \downarrow }^{\phantom{\dagger}}\hat{c}_{ \mathbf{p}+\mathbf{s}, \uparrow }^{\phantom{\dagger}}: \\
	&+1 \sum_{ \mathbf{q} \mathbf{p} \mathbf{r} \mathbf{s} } U( \mathbf{p}, \mathbf{r}, \mathbf{s} )  \alpha( \mathbf{p}+\mathbf{s}, \mathbf{q}, \mathbf{s} )   \langle \hat{n}_{ \mathbf{p}+\mathbf{s}, \uparrow } \rangle :\hat{c}_{ \mathbf{r}, \downarrow }^\dagger \hat{c}_{ \mathbf{q}, \downarrow }^\dagger \hat{c}_{ \mathbf{r}-\mathbf{s}, \downarrow }^{\phantom{\dagger}}\hat{c}_{ \mathbf{q}+\mathbf{s}, \downarrow }^{\phantom{\dagger}}: \\
	&+1 \sum_{ \mathbf{q} \mathbf{p} \mathbf{r} \mathbf{s} \mathbf{t} } U( \mathbf{q}, \mathbf{p}, \mathbf{r} )  \alpha( \mathbf{s}, \mathbf{t}, 0 )   \langle \hat{n}_{ \mathbf{t}, \downarrow } \rangle \langle \hat{n}_{ \mathbf{s}, \uparrow } \rangle :\hat{c}_{ \mathbf{q}, \uparrow }^\dagger \hat{c}_{ \mathbf{p}, \downarrow }^\dagger \hat{c}_{ \mathbf{p}+\mathbf{r}, \downarrow }^{\phantom{\dagger}}\hat{c}_{ \mathbf{q}-\mathbf{r}, \uparrow }^{\phantom{\dagger}}: \\
	&+1 \sum_{ \mathbf{q} \mathbf{p} \mathbf{r} \mathbf{s} } U( \mathbf{p}, \mathbf{r}, \mathbf{s} )  \alpha( \mathbf{q}, \mathbf{r}+\mathbf{s}, \mathbf{s} )   \langle \hat{n}_{ \mathbf{r}+\mathbf{s}, \downarrow } \rangle :\hat{c}_{ \mathbf{p}, \uparrow }^\dagger \hat{c}_{ \mathbf{q}, \uparrow }^\dagger \hat{c}_{ \mathbf{p}+\mathbf{s}, \uparrow }^{\phantom{\dagger}}\hat{c}_{ \mathbf{q}-\mathbf{s}, \uparrow }^{\phantom{\dagger}}: \\
	&+1 \sum_{ \mathbf{q} \mathbf{p} \mathbf{r} \mathbf{s} } U( \mathbf{p}, \mathbf{r}, \mathbf{s} )  \alpha( \mathbf{p}+\mathbf{s}, \mathbf{q}+\mathbf{r}, \mathbf{q} )   \langle \hat{n}_{ \mathbf{p}+\mathbf{s}, \uparrow } \rangle :\hat{c}_{ \mathbf{p}, \uparrow }^\dagger \hat{c}_{ \mathbf{q}+\mathbf{r}, \downarrow }^\dagger \hat{c}_{ \mathbf{r}-\mathbf{s}, \downarrow }^{\phantom{\dagger}}\hat{c}_{ \mathbf{p}+\mathbf{q}+\mathbf{s}, \uparrow }^{\phantom{\dagger}}: \\
	&-1 \sum_{ \mathbf{q} \mathbf{p} \mathbf{r} \mathbf{s} } U( \mathbf{p}, \mathbf{r}, \mathbf{s} )  \alpha( \mathbf{q}, \mathbf{r}, 0 )   \langle \hat{n}_{ \mathbf{q}, \uparrow } \rangle :\hat{c}_{ \mathbf{p}, \uparrow }^\dagger \hat{c}_{ \mathbf{r}, \downarrow }^\dagger \hat{c}_{ \mathbf{r}-\mathbf{s}, \downarrow }^{\phantom{\dagger}}\hat{c}_{ \mathbf{p}+\mathbf{s}, \uparrow }^{\phantom{\dagger}}: \\
	&-1 \sum_{ \mathbf{q} \mathbf{p} \mathbf{r} \mathbf{s} } U( \mathbf{r}, \mathbf{s}, 0 )  \alpha( \mathbf{q}, \mathbf{p}+\mathbf{s}, \mathbf{p} )   \langle \hat{n}_{ \mathbf{r}, \uparrow } \rangle :\hat{c}_{ \mathbf{q}, \uparrow }^\dagger \hat{c}_{ \mathbf{p}+\mathbf{s}, \downarrow }^\dagger \hat{c}_{ \mathbf{s}, \downarrow }^{\phantom{\dagger}}\hat{c}_{ \mathbf{p}+\mathbf{q}, \uparrow }^{\phantom{\dagger}}: \\
	&+1 \sum_{ \mathbf{q} \mathbf{p} \mathbf{r} \mathbf{s} } U( \mathbf{p}, \mathbf{r}, \mathbf{s} )  \alpha( \mathbf{p}-\mathbf{q}, \mathbf{q}+\mathbf{r}, \mathbf{q} )   \langle \hat{n}_{ \mathbf{p}, \uparrow } \rangle :\hat{c}_{ \mathbf{p}-\mathbf{q}, \uparrow }^\dagger \hat{c}_{ \mathbf{q}+\mathbf{r}, \downarrow }^\dagger \hat{c}_{ \mathbf{r}-\mathbf{s}, \downarrow }^{\phantom{\dagger}}\hat{c}_{ \mathbf{p}+\mathbf{s}, \uparrow }^{\phantom{\dagger}}: \\
	&+1 \sum_{ \mathbf{q} \mathbf{p} \mathbf{r} \mathbf{s} } U( \mathbf{p}-\mathbf{q}, \mathbf{q}+\mathbf{r}, \mathbf{q} )  \alpha( \mathbf{p}, \mathbf{r}, \mathbf{s} )   :\hat{c}_{ \mathbf{p}-\mathbf{q}, \uparrow }^\dagger \hat{c}_{ \mathbf{q}+\mathbf{r}, \downarrow }^\dagger \hat{c}_{ \mathbf{r}-\mathbf{s}, \downarrow }^{\phantom{\dagger}}\hat{c}_{ \mathbf{p}+\mathbf{s}, \uparrow }^{\phantom{\dagger}}: \\
	&-1 \sum_{ \mathbf{q} \mathbf{p} \mathbf{r} \mathbf{s} } U( \mathbf{p}, \mathbf{r}, \mathbf{s} )  \alpha( \mathbf{p}-\mathbf{q}, \mathbf{q}+\mathbf{r}, \mathbf{q} )   :\hat{c}_{ \mathbf{p}-\mathbf{q}, \uparrow }^\dagger \hat{c}_{ \mathbf{q}+\mathbf{r}, \downarrow }^\dagger \hat{c}_{ \mathbf{r}-\mathbf{s}, \downarrow }^{\phantom{\dagger}}\hat{c}_{ \mathbf{p}+\mathbf{s}, \uparrow }^{\phantom{\dagger}}: \\
	&-1 \sum_{ \mathbf{q} \mathbf{p} \mathbf{r} \mathbf{s} } U( \mathbf{q}, \mathbf{p}+\mathbf{s}, \mathbf{p} )  \alpha( \mathbf{r}, \mathbf{s}, 0 )   \langle \hat{n}_{ -\mathbf{s}, \downarrow } \rangle \langle \hat{n}_{ \mathbf{r}, \uparrow } \rangle :\hat{c}_{ \mathbf{q}, \uparrow }^\dagger \hat{c}_{ -\mathbf{p}+\mathbf{s}, \downarrow }^\dagger \hat{c}_{ \mathbf{s}, \downarrow }^{\phantom{\dagger}}\hat{c}_{ -\mathbf{p}+\mathbf{q}, \uparrow }^{\phantom{\dagger}}: \\
	&-1 \sum_{ \mathbf{q} \mathbf{p} \mathbf{r} \mathbf{s} } U( \mathbf{q}, \mathbf{p}, 0 )  \alpha( \mathbf{r}, \mathbf{s}, 0 )   \langle \hat{n}_{ \mathbf{r}, \uparrow } \rangle \langle \hat{n}_{ \mathbf{q}, \uparrow } \rangle :\hat{c}_{ \mathbf{s}, \downarrow }^\dagger \hat{c}_{ \mathbf{p}, \downarrow }^\dagger \hat{c}_{ \mathbf{s}, \downarrow }^{\phantom{\dagger}}\hat{c}_{ \mathbf{p}, \downarrow }^{\phantom{\dagger}}: \\
	&+1 \sum_{ \mathbf{q} \mathbf{p} \mathbf{r} \mathbf{s} } U( \mathbf{p}+\mathbf{q}, \mathbf{r}+\mathbf{s}, \mathbf{q} )  \alpha( \mathbf{p}, \mathbf{r}, \mathbf{s} )   \langle \hat{n}_{ \mathbf{r}+\mathbf{s}, \downarrow } \rangle \langle \hat{n}_{ \mathbf{p}, \uparrow } \rangle :\hat{c}_{ \mathbf{p}+\mathbf{q}, \uparrow }^\dagger \hat{c}_{ \mathbf{r}, \downarrow }^\dagger \hat{c}_{ \mathbf{q}+\mathbf{r}+\mathbf{s}, \downarrow }^{\phantom{\dagger}}\hat{c}_{ \mathbf{p}-\mathbf{s}, \uparrow }^{\phantom{\dagger}}: \\
	&-1 \sum_{ \mathbf{q} \mathbf{p} \mathbf{r} \mathbf{s} } U( \mathbf{p}, \mathbf{q}, 0 )  \alpha( \mathbf{p}, \mathbf{r}, \mathbf{s} )   \langle \hat{n}_{ \mathbf{q}, \downarrow } \rangle \langle \hat{n}_{ \mathbf{p}, \uparrow } \rangle :\hat{c}_{ \mathbf{p}, \uparrow }^\dagger \hat{c}_{ \mathbf{r}, \downarrow }^\dagger \hat{c}_{ \mathbf{r}+\mathbf{s}, \downarrow }^{\phantom{\dagger}}\hat{c}_{ \mathbf{p}-\mathbf{s}, \uparrow }^{\phantom{\dagger}}: \\
	&-1 \sum_{ \mathbf{q} \mathbf{p} \mathbf{r} \mathbf{s} } U( \mathbf{p}+\mathbf{r}, \mathbf{q}, \mathbf{p} )  \alpha( \mathbf{r}, \mathbf{s}, 0 )   \langle \hat{n}_{ \mathbf{s}, \downarrow } \rangle \langle \hat{n}_{ \mathbf{r}, \uparrow } \rangle :\hat{c}_{ \mathbf{p}+\mathbf{r}, \uparrow }^\dagger \hat{c}_{ \mathbf{q}, \downarrow }^\dagger \hat{c}_{ \mathbf{p}+\mathbf{q}, \downarrow }^{\phantom{\dagger}}\hat{c}_{ \mathbf{r}, \uparrow }^{\phantom{\dagger}}: \\
	&+1 \sum_{ \mathbf{q} \mathbf{p} \mathbf{r} \mathbf{s} } U( \mathbf{p}+\mathbf{q}, \mathbf{q}+\mathbf{r}, \mathbf{q} )  \alpha( \mathbf{p}, \mathbf{r}, \mathbf{s} )   \langle \hat{n}_{ -\mathbf{r}, \downarrow } \rangle \langle \hat{n}_{ \mathbf{p}, \uparrow } \rangle :\hat{c}_{ \mathbf{p}+\mathbf{q}, \uparrow }^\dagger \hat{c}_{ -\mathbf{q}+\mathbf{r}, \downarrow }^\dagger \hat{c}_{ \mathbf{r}+\mathbf{s}, \downarrow }^{\phantom{\dagger}}\hat{c}_{ \mathbf{p}-\mathbf{s}, \uparrow }^{\phantom{\dagger}}: \\
	&+1 \sum_{ \mathbf{q} \mathbf{p} \mathbf{r} \mathbf{s} } U( \mathbf{p}+\mathbf{s}, \mathbf{q}, \mathbf{s} )  \alpha( \mathbf{p}, \mathbf{r}, \mathbf{s} )   \langle \hat{n}_{ \mathbf{p}, \uparrow } \rangle \langle \hat{n}_{ \mathbf{p}+\mathbf{s}, \uparrow } \rangle :\hat{c}_{ \mathbf{r}, \downarrow }^\dagger \hat{c}_{ \mathbf{q}, \downarrow }^\dagger \hat{c}_{ \mathbf{r}+\mathbf{s}, \downarrow }^{\phantom{\dagger}}\hat{c}_{ \mathbf{q}-\mathbf{s}, \downarrow }^{\phantom{\dagger}}: \\
	&-1 \sum_{ \mathbf{q} \mathbf{p} \mathbf{r} \mathbf{s} } U( \mathbf{q}, \mathbf{p}, 0 )  \alpha( \mathbf{r}, \mathbf{s}, 0 )   \langle \hat{n}_{ \mathbf{s}, \downarrow } \rangle \langle \hat{n}_{ \mathbf{p}, \downarrow } \rangle :\hat{c}_{ \mathbf{r}, \uparrow }^\dagger \hat{c}_{ \mathbf{q}, \uparrow }^\dagger \hat{c}_{ \mathbf{r}, \uparrow }^{\phantom{\dagger}}\hat{c}_{ \mathbf{q}, \uparrow }^{\phantom{\dagger}}: \\
	&+1 \sum_{ \mathbf{q} \mathbf{p} \mathbf{r} \mathbf{s} } U( \mathbf{q}, \mathbf{r}+\mathbf{s}, \mathbf{s} )  \alpha( \mathbf{p}, \mathbf{r}, \mathbf{s} )   \langle \hat{n}_{ \mathbf{r}, \downarrow } \rangle \langle \hat{n}_{ \mathbf{r}+\mathbf{s}, \downarrow } \rangle :\hat{c}_{ \mathbf{p}, \uparrow }^\dagger \hat{c}_{ \mathbf{q}, \uparrow }^\dagger \hat{c}_{ \mathbf{p}-\mathbf{s}, \uparrow }^{\phantom{\dagger}}\hat{c}_{ \mathbf{q}+\mathbf{s}, \uparrow }^{\phantom{\dagger}}: \\
	&+1 \sum_{ \mathbf{q} \mathbf{p} \mathbf{r} \mathbf{s} } U( \mathbf{p}-\mathbf{s}, \mathbf{r}+\mathbf{s}, \mathbf{q} )  \alpha( \mathbf{p}, \mathbf{r}, \mathbf{s} )   \langle \hat{n}_{ \mathbf{r}+\mathbf{s}, \downarrow } \rangle \langle \hat{n}_{ \mathbf{p}-\mathbf{s}, \uparrow } \rangle :\hat{c}_{ \mathbf{p}, \uparrow }^\dagger \hat{c}_{ \mathbf{r}, \downarrow }^\dagger \hat{c}_{ \mathbf{q}+\mathbf{r}+\mathbf{s}, \downarrow }^{\phantom{\dagger}}\hat{c}_{ \mathbf{p}-\mathbf{q}-\mathbf{s}, \uparrow }^{\phantom{\dagger}}: \\
	&-1 \sum_{ \mathbf{q} \mathbf{p} \mathbf{r} \mathbf{s} } U( \mathbf{p}+\mathbf{s}, \mathbf{q}, 0 )  \alpha( \mathbf{p}, \mathbf{r}, \mathbf{s} )   \langle \hat{n}_{ \mathbf{q}, \downarrow } \rangle \langle \hat{n}_{ \mathbf{p}+\mathbf{s}, \uparrow } \rangle :\hat{c}_{ \mathbf{p}, \uparrow }^\dagger \hat{c}_{ \mathbf{r}, \downarrow }^\dagger \hat{c}_{ \mathbf{r}+\mathbf{s}, \downarrow }^{\phantom{\dagger}}\hat{c}_{ \mathbf{p}-\mathbf{s}, \uparrow }^{\phantom{\dagger}}: \\
	&-1 \sum_{ \mathbf{q} \mathbf{p} \mathbf{r} \mathbf{s} } U( \mathbf{r}, \mathbf{q}, \mathbf{p} )  \alpha( \mathbf{r}, \mathbf{s}, 0 )   \langle \hat{n}_{ \mathbf{s}, \downarrow } \rangle \langle \hat{n}_{ \mathbf{r}, \uparrow } \rangle :\hat{c}_{ \mathbf{r}, \uparrow }^\dagger \hat{c}_{ \mathbf{q}, \downarrow }^\dagger \hat{c}_{ \mathbf{p}+\mathbf{q}, \downarrow }^{\phantom{\dagger}}\hat{c}_{ -\mathbf{p}+\mathbf{r}, \uparrow }^{\phantom{\dagger}}: \\
	&+1 \sum_{ \mathbf{q} \mathbf{p} \mathbf{r} \mathbf{s} } U( \mathbf{p}+\mathbf{s}, \mathbf{q}+\mathbf{r}, \mathbf{q} )  \alpha( \mathbf{p}, \mathbf{r}, \mathbf{s} )   \langle \hat{n}_{ -\mathbf{r}, \downarrow } \rangle \langle \hat{n}_{ \mathbf{p}+\mathbf{s}, \uparrow } \rangle :\hat{c}_{ \mathbf{p}, \uparrow }^\dagger \hat{c}_{ -\mathbf{q}+\mathbf{r}, \downarrow }^\dagger \hat{c}_{ \mathbf{r}+\mathbf{s}, \downarrow }^{\phantom{\dagger}}\hat{c}_{ \mathbf{p}-\mathbf{q}-\mathbf{s}, \uparrow }^{\phantom{\dagger}}: \\
	&-1 \sum_{ \mathbf{q} \mathbf{p} \mathbf{r} \mathbf{s} } U( \mathbf{q}, \mathbf{r}+\mathbf{s}, 0 )  \alpha( \mathbf{p}, \mathbf{r}, \mathbf{s} )   \langle \hat{n}_{ \mathbf{r}+\mathbf{s}, \downarrow } \rangle \langle \hat{n}_{ \mathbf{q}, \uparrow } \rangle :\hat{c}_{ \mathbf{p}, \uparrow }^\dagger \hat{c}_{ \mathbf{r}, \downarrow }^\dagger \hat{c}_{ \mathbf{r}+\mathbf{s}, \downarrow }^{\phantom{\dagger}}\hat{c}_{ \mathbf{p}-\mathbf{s}, \uparrow }^{\phantom{\dagger}}: \\
	&+1 \sum_{ \mathbf{q} \mathbf{p} \mathbf{r} \mathbf{s} \mathbf{t} } U( \mathbf{q}, \mathbf{p}, 0 )  \alpha( \mathbf{r}, \mathbf{s}, \mathbf{t} )   \langle \hat{n}_{ \mathbf{q}, \uparrow } \rangle \langle \hat{n}_{ \mathbf{p}, \downarrow } \rangle :\hat{c}_{ \mathbf{r}, \uparrow }^\dagger \hat{c}_{ \mathbf{s}, \downarrow }^\dagger \hat{c}_{ \mathbf{s}+\mathbf{t}, \downarrow }^{\phantom{\dagger}}\hat{c}_{ \mathbf{r}-\mathbf{t}, \uparrow }^{\phantom{\dagger}}: \\
	&+1 \sum_{ \mathbf{q} \mathbf{p} \mathbf{r} \mathbf{s} } U( \mathbf{q}, \mathbf{p}, 0 )  \alpha( \mathbf{r}, \mathbf{s}, 0 )   \langle \hat{n}_{ \mathbf{s}, \downarrow } \rangle \langle \hat{n}_{ \mathbf{q}, \uparrow } \rangle :\hat{c}_{ \mathbf{r}, \uparrow }^\dagger \hat{c}_{ \mathbf{p}, \downarrow }^\dagger \hat{c}_{ \mathbf{p}, \downarrow }^{\phantom{\dagger}}\hat{c}_{ \mathbf{r}, \uparrow }^{\phantom{\dagger}}: \\
	&+1 \sum_{ \mathbf{q} \mathbf{p} \mathbf{r} \mathbf{s} } U( \mathbf{r}, \mathbf{s}, 0 )  \alpha( \mathbf{q}, \mathbf{p}, 0 )   \langle \hat{n}_{ \mathbf{s}, \downarrow } \rangle \langle \hat{n}_{ \mathbf{p}, \downarrow } \rangle :\hat{c}_{ \mathbf{r}, \uparrow }^\dagger \hat{c}_{ \mathbf{q}, \uparrow }^\dagger \hat{c}_{ \mathbf{r}, \uparrow }^{\phantom{\dagger}}\hat{c}_{ \mathbf{q}, \uparrow }^{\phantom{\dagger}}: \\
	&-1 \sum_{ \mathbf{q} \mathbf{p} \mathbf{r} \mathbf{s} } U( \mathbf{p}, \mathbf{r}, \mathbf{s} )  \alpha( \mathbf{q}, \mathbf{r}+\mathbf{s}, \mathbf{s} )   \langle \hat{n}_{ \mathbf{r}, \downarrow } \rangle \langle \hat{n}_{ \mathbf{r}+\mathbf{s}, \downarrow } \rangle :\hat{c}_{ \mathbf{p}, \uparrow }^\dagger \hat{c}_{ \mathbf{q}, \uparrow }^\dagger \hat{c}_{ \mathbf{p}-\mathbf{s}, \uparrow }^{\phantom{\dagger}}\hat{c}_{ \mathbf{q}+\mathbf{s}, \uparrow }^{\phantom{\dagger}}: \\
	&-1 \sum_{ \mathbf{q} \mathbf{p} \mathbf{r} \mathbf{s} } U( \mathbf{p}, \mathbf{r}, \mathbf{s} )  \alpha( \mathbf{p}-\mathbf{s}, \mathbf{r}+\mathbf{s}, \mathbf{q} )   \langle \hat{n}_{ \mathbf{r}+\mathbf{s}, \downarrow } \rangle \langle \hat{n}_{ \mathbf{p}-\mathbf{s}, \uparrow } \rangle :\hat{c}_{ \mathbf{p}, \uparrow }^\dagger \hat{c}_{ \mathbf{r}, \downarrow }^\dagger \hat{c}_{ \mathbf{q}+\mathbf{r}+\mathbf{s}, \downarrow }^{\phantom{\dagger}}\hat{c}_{ \mathbf{p}-\mathbf{q}-\mathbf{s}, \uparrow }^{\phantom{\dagger}}: \\
	&+1 \sum_{ \mathbf{q} \mathbf{p} \mathbf{r} \mathbf{s} } U( \mathbf{p}, \mathbf{r}, \mathbf{s} )  \alpha( \mathbf{p}+\mathbf{s}, \mathbf{q}, 0 )   \langle \hat{n}_{ \mathbf{q}, \downarrow } \rangle \langle \hat{n}_{ \mathbf{p}+\mathbf{s}, \uparrow } \rangle :\hat{c}_{ \mathbf{p}, \uparrow }^\dagger \hat{c}_{ \mathbf{r}, \downarrow }^\dagger \hat{c}_{ \mathbf{r}+\mathbf{s}, \downarrow }^{\phantom{\dagger}}\hat{c}_{ \mathbf{p}-\mathbf{s}, \uparrow }^{\phantom{\dagger}}: \\
	&+1 \sum_{ \mathbf{q} \mathbf{p} \mathbf{r} \mathbf{s} } U( \mathbf{r}, \mathbf{s}, 0 )  \alpha( \mathbf{r}, \mathbf{q}, \mathbf{p} )   \langle \hat{n}_{ \mathbf{s}, \downarrow } \rangle \langle \hat{n}_{ \mathbf{r}, \uparrow } \rangle :\hat{c}_{ \mathbf{r}, \uparrow }^\dagger \hat{c}_{ \mathbf{q}, \downarrow }^\dagger \hat{c}_{ \mathbf{p}+\mathbf{q}, \downarrow }^{\phantom{\dagger}}\hat{c}_{ -\mathbf{p}+\mathbf{r}, \uparrow }^{\phantom{\dagger}}: \\
	&-1 \sum_{ \mathbf{q} \mathbf{p} \mathbf{r} \mathbf{s} } U( \mathbf{p}, \mathbf{r}, \mathbf{s} )  \alpha( \mathbf{p}+\mathbf{s}, \mathbf{q}+\mathbf{r}, \mathbf{q} )   \langle \hat{n}_{ -\mathbf{r}, \downarrow } \rangle \langle \hat{n}_{ \mathbf{p}+\mathbf{s}, \uparrow } \rangle :\hat{c}_{ \mathbf{p}, \uparrow }^\dagger \hat{c}_{ -\mathbf{q}+\mathbf{r}, \downarrow }^\dagger \hat{c}_{ \mathbf{r}+\mathbf{s}, \downarrow }^{\phantom{\dagger}}\hat{c}_{ \mathbf{p}-\mathbf{q}-\mathbf{s}, \uparrow }^{\phantom{\dagger}}: \\
	&+1 \sum_{ \mathbf{q} \mathbf{p} \mathbf{r} \mathbf{s} } U( \mathbf{p}, \mathbf{r}, \mathbf{s} )  \alpha( \mathbf{q}, \mathbf{r}+\mathbf{s}, 0 )   \langle \hat{n}_{ \mathbf{r}+\mathbf{s}, \downarrow } \rangle \langle \hat{n}_{ \mathbf{q}, \uparrow } \rangle :\hat{c}_{ \mathbf{p}, \uparrow }^\dagger \hat{c}_{ \mathbf{r}, \downarrow }^\dagger \hat{c}_{ \mathbf{r}+\mathbf{s}, \downarrow }^{\phantom{\dagger}}\hat{c}_{ \mathbf{p}-\mathbf{s}, \uparrow }^{\phantom{\dagger}}: \\
	&-1 \sum_{ \mathbf{q} \mathbf{p} \mathbf{r} \mathbf{s} \mathbf{t} } U( \mathbf{r}, \mathbf{s}, \mathbf{t} )  \alpha( \mathbf{q}, \mathbf{p}, 0 )   \langle \hat{n}_{ \mathbf{q}, \uparrow } \rangle \langle \hat{n}_{ \mathbf{p}, \downarrow } \rangle :\hat{c}_{ \mathbf{r}, \uparrow }^\dagger \hat{c}_{ \mathbf{s}, \downarrow }^\dagger \hat{c}_{ \mathbf{s}+\mathbf{t}, \downarrow }^{\phantom{\dagger}}\hat{c}_{ \mathbf{r}-\mathbf{t}, \uparrow }^{\phantom{\dagger}}: \\
	&-1 \sum_{ \mathbf{q} \mathbf{p} \mathbf{r} \mathbf{s} } U( \mathbf{r}, \mathbf{s}, 0 )  \alpha( \mathbf{q}, \mathbf{p}, 0 )   \langle \hat{n}_{ \mathbf{s}, \downarrow } \rangle \langle \hat{n}_{ \mathbf{q}, \uparrow } \rangle :\hat{c}_{ \mathbf{r}, \uparrow }^\dagger \hat{c}_{ \mathbf{p}, \downarrow }^\dagger \hat{c}_{ \mathbf{p}, \downarrow }^{\phantom{\dagger}}\hat{c}_{ \mathbf{r}, \uparrow }^{\phantom{\dagger}}: \\
	&+1 \sum_{ \mathbf{q} \mathbf{p} \mathbf{r} \mathbf{s} } U( \mathbf{p}, \mathbf{r}, \mathbf{s} )  \alpha( \mathbf{q}, \mathbf{r}, 0 )   \langle \hat{n}_{ \mathbf{r}, \downarrow } \rangle \langle \hat{n}_{ \mathbf{q}, \uparrow } \rangle :\hat{c}_{ \mathbf{p}, \uparrow }^\dagger \hat{c}_{ \mathbf{r}, \downarrow }^\dagger \hat{c}_{ \mathbf{r}+\mathbf{s}, \downarrow }^{\phantom{\dagger}}\hat{c}_{ \mathbf{p}-\mathbf{s}, \uparrow }^{\phantom{\dagger}}: \\
	&+1 \sum_{ \mathbf{q} \mathbf{p} \mathbf{r} \mathbf{s} } U( \mathbf{r}, \mathbf{s}, 0 )  \alpha( \mathbf{q}, \mathbf{s}, \mathbf{p} )   \langle \hat{n}_{ \mathbf{s}, \downarrow } \rangle \langle \hat{n}_{ \mathbf{r}, \uparrow } \rangle :\hat{c}_{ \mathbf{q}, \uparrow }^\dagger \hat{c}_{ \mathbf{s}, \downarrow }^\dagger \hat{c}_{ \mathbf{p}+\mathbf{s}, \downarrow }^{\phantom{\dagger}}\hat{c}_{ -\mathbf{p}+\mathbf{q}, \uparrow }^{\phantom{\dagger}}: \\
	&-1 \sum_{ \mathbf{q} \mathbf{p} \mathbf{r} \mathbf{s} } U( \mathbf{r}, \mathbf{s}, 0 )  \alpha( \mathbf{q}, \mathbf{p}, 0 )   \langle \hat{n}_{ \mathbf{r}, \uparrow } \rangle \langle \hat{n}_{ \mathbf{p}, \downarrow } \rangle :\hat{c}_{ \mathbf{q}, \uparrow }^\dagger \hat{c}_{ \mathbf{s}, \downarrow }^\dagger \hat{c}_{ \mathbf{s}, \downarrow }^{\phantom{\dagger}}\hat{c}_{ \mathbf{q}, \uparrow }^{\phantom{\dagger}}: \\
	&-1 \sum_{ \mathbf{q} \mathbf{p} \mathbf{r} \mathbf{s} \mathbf{t} } U( \mathbf{s}, \mathbf{t}, 0 )  \alpha( \mathbf{q}, \mathbf{p}, \mathbf{r} )   \langle \hat{n}_{ \mathbf{t}, \downarrow } \rangle \langle \hat{n}_{ \mathbf{s}, \uparrow } \rangle :\hat{c}_{ \mathbf{q}, \uparrow }^\dagger \hat{c}_{ \mathbf{p}, \downarrow }^\dagger \hat{c}_{ \mathbf{p}+\mathbf{r}, \downarrow }^{\phantom{\dagger}}\hat{c}_{ \mathbf{q}-\mathbf{r}, \uparrow }^{\phantom{\dagger}}: \\
	&+1 \sum_{ \mathbf{q} \mathbf{p} \mathbf{r} \mathbf{s} } U( \mathbf{r}, \mathbf{s}, 0 )  \alpha( \mathbf{q}, \mathbf{p}+\mathbf{s}, \mathbf{p} )   \langle \hat{n}_{ -\mathbf{s}, \downarrow } \rangle \langle \hat{n}_{ \mathbf{r}, \uparrow } \rangle :\hat{c}_{ \mathbf{q}, \uparrow }^\dagger \hat{c}_{ -\mathbf{p}+\mathbf{s}, \downarrow }^\dagger \hat{c}_{ \mathbf{s}, \downarrow }^{\phantom{\dagger}}\hat{c}_{ -\mathbf{p}+\mathbf{q}, \uparrow }^{\phantom{\dagger}}: \\
	&+1 \sum_{ \mathbf{q} \mathbf{p} \mathbf{r} \mathbf{s} } U( \mathbf{r}, \mathbf{s}, 0 )  \alpha( \mathbf{q}, \mathbf{p}, 0 )   \langle \hat{n}_{ \mathbf{r}, \uparrow } \rangle \langle \hat{n}_{ \mathbf{q}, \uparrow } \rangle :\hat{c}_{ \mathbf{s}, \downarrow }^\dagger \hat{c}_{ \mathbf{p}, \downarrow }^\dagger \hat{c}_{ \mathbf{s}, \downarrow }^{\phantom{\dagger}}\hat{c}_{ \mathbf{p}, \downarrow }^{\phantom{\dagger}}: \\
	&-1 \sum_{ \mathbf{q} \mathbf{p} \mathbf{r} \mathbf{s} } U( \mathbf{p}, \mathbf{r}, \mathbf{s} )  \alpha( \mathbf{p}+\mathbf{q}, \mathbf{r}+\mathbf{s}, \mathbf{q} )   \langle \hat{n}_{ \mathbf{r}+\mathbf{s}, \downarrow } \rangle \langle \hat{n}_{ \mathbf{p}, \uparrow } \rangle :\hat{c}_{ \mathbf{p}+\mathbf{q}, \uparrow }^\dagger \hat{c}_{ \mathbf{r}, \downarrow }^\dagger \hat{c}_{ \mathbf{q}+\mathbf{r}+\mathbf{s}, \downarrow }^{\phantom{\dagger}}\hat{c}_{ \mathbf{p}-\mathbf{s}, \uparrow }^{\phantom{\dagger}}: \\
	&+1 \sum_{ \mathbf{q} \mathbf{p} \mathbf{r} \mathbf{s} } U( \mathbf{p}, \mathbf{r}, \mathbf{s} )  \alpha( \mathbf{p}, \mathbf{q}, 0 )   \langle \hat{n}_{ \mathbf{q}, \downarrow } \rangle \langle \hat{n}_{ \mathbf{p}, \uparrow } \rangle :\hat{c}_{ \mathbf{p}, \uparrow }^\dagger \hat{c}_{ \mathbf{r}, \downarrow }^\dagger \hat{c}_{ \mathbf{r}+\mathbf{s}, \downarrow }^{\phantom{\dagger}}\hat{c}_{ \mathbf{p}-\mathbf{s}, \uparrow }^{\phantom{\dagger}}: \\
	&+1 \sum_{ \mathbf{q} \mathbf{p} \mathbf{r} \mathbf{s} } U( \mathbf{r}, \mathbf{s}, 0 )  \alpha( \mathbf{p}+\mathbf{r}, \mathbf{q}, \mathbf{p} )   \langle \hat{n}_{ \mathbf{s}, \downarrow } \rangle \langle \hat{n}_{ \mathbf{r}, \uparrow } \rangle :\hat{c}_{ \mathbf{p}+\mathbf{r}, \uparrow }^\dagger \hat{c}_{ \mathbf{q}, \downarrow }^\dagger \hat{c}_{ \mathbf{p}+\mathbf{q}, \downarrow }^{\phantom{\dagger}}\hat{c}_{ \mathbf{r}, \uparrow }^{\phantom{\dagger}}: \\
	&-1 \sum_{ \mathbf{q} \mathbf{p} \mathbf{r} \mathbf{s} } U( \mathbf{p}, \mathbf{r}, \mathbf{s} )  \alpha( \mathbf{p}+\mathbf{q}, \mathbf{q}+\mathbf{r}, \mathbf{q} )   \langle \hat{n}_{ -\mathbf{r}, \downarrow } \rangle \langle \hat{n}_{ \mathbf{p}, \uparrow } \rangle :\hat{c}_{ \mathbf{p}+\mathbf{q}, \uparrow }^\dagger \hat{c}_{ -\mathbf{q}+\mathbf{r}, \downarrow }^\dagger \hat{c}_{ \mathbf{r}+\mathbf{s}, \downarrow }^{\phantom{\dagger}}\hat{c}_{ \mathbf{p}-\mathbf{s}, \uparrow }^{\phantom{\dagger}}: \\
	&-1 \sum_{ \mathbf{q} \mathbf{p} \mathbf{r} \mathbf{s} } U( \mathbf{p}, \mathbf{r}, \mathbf{s} )  \alpha( \mathbf{p}+\mathbf{s}, \mathbf{q}, \mathbf{s} )   \langle \hat{n}_{ \mathbf{p}, \uparrow } \rangle \langle \hat{n}_{ \mathbf{p}+\mathbf{s}, \uparrow } \rangle :\hat{c}_{ \mathbf{r}, \downarrow }^\dagger \hat{c}_{ \mathbf{q}, \downarrow }^\dagger \hat{c}_{ \mathbf{r}+\mathbf{s}, \downarrow }^{\phantom{\dagger}}\hat{c}_{ \mathbf{q}-\mathbf{s}, \downarrow }^{\phantom{\dagger}}: \\
	&-1 \sum_{ \mathbf{q} \mathbf{p} \mathbf{r} \mathbf{s} } U( \mathbf{q}, \mathbf{r}, 0 )  \alpha( \mathbf{p}, \mathbf{r}, \mathbf{s} )   \langle \hat{n}_{ \mathbf{r}, \downarrow } \rangle \langle \hat{n}_{ \mathbf{q}, \uparrow } \rangle :\hat{c}_{ \mathbf{p}, \uparrow }^\dagger \hat{c}_{ \mathbf{r}, \downarrow }^\dagger \hat{c}_{ \mathbf{r}+\mathbf{s}, \downarrow }^{\phantom{\dagger}}\hat{c}_{ \mathbf{p}-\mathbf{s}, \uparrow }^{\phantom{\dagger}}: \\
	&-1 \sum_{ \mathbf{q} \mathbf{p} \mathbf{r} \mathbf{s} } U( \mathbf{p}+\mathbf{q}, \mathbf{r}+\mathbf{s}, \mathbf{q} )  \alpha( \mathbf{p}, \mathbf{r}, \mathbf{s} )   \langle \hat{n}_{ \mathbf{r}+\mathbf{s}, \downarrow } \rangle :\hat{c}_{ \mathbf{p}-\mathbf{q}, \uparrow }^\dagger \hat{c}_{ \mathbf{r}, \downarrow }^\dagger \hat{c}_{ -\mathbf{q}+\mathbf{r}-\mathbf{s}, \downarrow }^{\phantom{\dagger}}\hat{c}_{ \mathbf{p}+\mathbf{s}, \uparrow }^{\phantom{\dagger}}: \\
	&+1 \sum_{ \mathbf{q} \mathbf{p} \mathbf{r} \mathbf{s} } U( \mathbf{p}, \mathbf{q}, 0 )  \alpha( \mathbf{p}, \mathbf{r}, \mathbf{s} )   \langle \hat{n}_{ \mathbf{q}, \downarrow } \rangle :\hat{c}_{ \mathbf{p}, \uparrow }^\dagger \hat{c}_{ \mathbf{r}, \downarrow }^\dagger \hat{c}_{ \mathbf{r}-\mathbf{s}, \downarrow }^{\phantom{\dagger}}\hat{c}_{ \mathbf{p}+\mathbf{s}, \uparrow }^{\phantom{\dagger}}: \\
	&+1 \sum_{ \mathbf{q} \mathbf{p} \mathbf{r} \mathbf{s} } U( \mathbf{p}+\mathbf{r}, \mathbf{q}, \mathbf{p} )  \alpha( \mathbf{r}, \mathbf{s}, 0 )   \langle \hat{n}_{ \mathbf{s}, \downarrow } \rangle :\hat{c}_{ -\mathbf{p}+\mathbf{r}, \uparrow }^\dagger \hat{c}_{ \mathbf{q}, \downarrow }^\dagger \hat{c}_{ -\mathbf{p}+\mathbf{q}, \downarrow }^{\phantom{\dagger}}\hat{c}_{ \mathbf{r}, \uparrow }^{\phantom{\dagger}}: \\
	&+2 \sum_{ \mathbf{q} \mathbf{p} \mathbf{r} } U( \mathbf{p}, \mathbf{r}, 0 )  \alpha( \mathbf{q}, \mathbf{r}, 0 )   \langle \hat{n}_{ \mathbf{r}, \downarrow } \rangle \langle \hat{n}_{ \mathbf{q}, \uparrow } \rangle \langle \hat{n}_{ \mathbf{p}, \uparrow } \rangle :\hat{c}_{ \mathbf{r}, \downarrow }^\dagger \hat{c}_{ \mathbf{r}, \downarrow }^{\phantom{\dagger}}: \\
	&-1 \sum_{ \mathbf{q} \mathbf{p} \mathbf{r} \mathbf{s} } U( \mathbf{r}, \mathbf{s}, 0 )  \alpha( \mathbf{q}, \mathbf{p}, 0 )   \langle \hat{n}_{ \mathbf{r}, \uparrow } \rangle \langle \hat{n}_{ \mathbf{q}, \uparrow } \rangle \langle \hat{n}_{ \mathbf{p}, \downarrow } \rangle :\hat{c}_{ \mathbf{s}, \downarrow }^\dagger \hat{c}_{ \mathbf{s}, \downarrow }^{\phantom{\dagger}}: \\
	&-1 \sum_{ \mathbf{q} \mathbf{p} \mathbf{r} \mathbf{s} } U( \mathbf{r}, \mathbf{s}, 0 )  \alpha( \mathbf{q}, \mathbf{p}, 0 )   \langle \hat{n}_{ \mathbf{s}, \downarrow } \rangle \langle \hat{n}_{ \mathbf{r}, \uparrow } \rangle \langle \hat{n}_{ \mathbf{q}, \uparrow } \rangle :\hat{c}_{ \mathbf{p}, \downarrow }^\dagger \hat{c}_{ \mathbf{p}, \downarrow }^{\phantom{\dagger}}: \\
	&-1 \sum_{ \mathbf{q} \mathbf{p} \mathbf{r} } U( \mathbf{q}, \mathbf{p}, \mathbf{r} )  \alpha( \mathbf{q}-\mathbf{r}, \mathbf{p}+\mathbf{r}, \mathbf{r} )   \langle \hat{n}_{ \mathbf{q}, \uparrow } \rangle \langle \hat{n}_{ \mathbf{p}, \downarrow } \rangle \langle \hat{n}_{ \mathbf{p}+\mathbf{r}, \downarrow } \rangle :\hat{c}_{ \mathbf{q}-\mathbf{r}, \uparrow }^\dagger \hat{c}_{ \mathbf{q}-\mathbf{r}, \uparrow }^{\phantom{\dagger}}: \\
	&+1 \sum_{ \mathbf{q} \mathbf{p} \mathbf{r} } U( \mathbf{q}, \mathbf{p}, \mathbf{r} )  \alpha( \mathbf{q}-\mathbf{r}, \mathbf{p}+\mathbf{r}, \mathbf{r} )   \langle \hat{n}_{ \mathbf{q}-\mathbf{r}, \uparrow } \rangle \langle \hat{n}_{ \mathbf{p}+\mathbf{r}, \downarrow } \rangle :\hat{c}_{ \mathbf{q}, \uparrow }^\dagger \hat{c}_{ \mathbf{q}, \uparrow }^{\phantom{\dagger}}: \\
	&-1 \sum_{ \mathbf{q} \mathbf{p} \mathbf{r} } U( \mathbf{p}, \mathbf{r}, 0 )  \alpha( \mathbf{q}, \mathbf{r}, 0 )   \langle \hat{n}_{ \mathbf{r}, \downarrow } \rangle \langle \hat{n}_{ \mathbf{q}, \uparrow } \rangle :\hat{c}_{ \mathbf{p}, \uparrow }^\dagger \hat{c}_{ \mathbf{p}, \uparrow }^{\phantom{\dagger}}: \\
	&-1 \sum_{ \mathbf{q} \mathbf{p} \mathbf{r} } U( \mathbf{p}, \mathbf{r}, 0 )  \alpha( \mathbf{q}, \mathbf{r}, 0 )   \langle \hat{n}_{ \mathbf{r}, \downarrow } \rangle \langle \hat{n}_{ \mathbf{p}, \uparrow } \rangle :\hat{c}_{ \mathbf{q}, \uparrow }^\dagger \hat{c}_{ \mathbf{q}, \uparrow }^{\phantom{\dagger}}: \\
	&-1 \sum_{ \mathbf{q} \mathbf{p} \mathbf{r} } U( \mathbf{p}, \mathbf{r}, 0 )  \alpha( \mathbf{q}, \mathbf{r}, 0 )   \langle \hat{n}_{ \mathbf{q}, \uparrow } \rangle \langle \hat{n}_{ \mathbf{p}, \uparrow } \rangle :\hat{c}_{ \mathbf{r}, \downarrow }^\dagger \hat{c}_{ \mathbf{r}, \downarrow }^{\phantom{\dagger}}: \\
	&+1 \sum_{ \mathbf{q} \mathbf{p} \mathbf{r} } U( \mathbf{q}, \mathbf{p}, \mathbf{r} )  \alpha( \mathbf{q}-\mathbf{r}, \mathbf{p}+\mathbf{r}, \mathbf{r} )   \langle \hat{n}_{ \mathbf{q}, \uparrow } \rangle \langle \hat{n}_{ \mathbf{p}+\mathbf{r}, \downarrow } \rangle :\hat{c}_{ \mathbf{q}+\mathbf{r}, \uparrow }^\dagger \hat{c}_{ \mathbf{q}+\mathbf{r}, \uparrow }^{\phantom{\dagger}}: \\
	&+1 \sum_{ \mathbf{q} \mathbf{p} \mathbf{r} } U( \mathbf{q}, \mathbf{p}, \mathbf{r} )  \alpha( \mathbf{q}-\mathbf{r}, \mathbf{p}+\mathbf{r}, \mathbf{r} )   \langle \hat{n}_{ \mathbf{q}, \uparrow } \rangle \langle \hat{n}_{ \mathbf{q}-\mathbf{r}, \uparrow } \rangle :\hat{c}_{ \mathbf{p}-\mathbf{r}, \downarrow }^\dagger \hat{c}_{ \mathbf{p}-\mathbf{r}, \downarrow }^{\phantom{\dagger}}: \\
	&-1 \sum_{ \mathbf{q} \mathbf{p} \mathbf{r} } U( \mathbf{q}, \mathbf{p}, \mathbf{r} )  \alpha( \mathbf{q}-\mathbf{r}, \mathbf{p}+\mathbf{r}, \mathbf{r} )   \langle \hat{n}_{ \mathbf{q}, \uparrow } \rangle \langle \hat{n}_{ \mathbf{q}-\mathbf{r}, \uparrow } \rangle \langle \hat{n}_{ \mathbf{p}+\mathbf{r}, \downarrow } \rangle :\hat{c}_{ \mathbf{p}, \downarrow }^\dagger \hat{c}_{ \mathbf{p}, \downarrow }^{\phantom{\dagger}}: \\
	&+1 \sum_{ \mathbf{q} \mathbf{p} \mathbf{r} } U( \mathbf{q}-\mathbf{r}, \mathbf{p}+\mathbf{r}, \mathbf{r} )  \alpha( \mathbf{q}, \mathbf{p}, \mathbf{r} )   \langle \hat{n}_{ \mathbf{p}+\mathbf{r}, \downarrow } \rangle :\hat{c}_{ \mathbf{q}+\mathbf{r}, \uparrow }^\dagger \hat{c}_{ \mathbf{q}+\mathbf{r}, \uparrow }^{\phantom{\dagger}}: \\
	&+1 \sum_{ \mathbf{q} \mathbf{p} \mathbf{r} } U( \mathbf{q}-\mathbf{r}, \mathbf{p}+\mathbf{r}, \mathbf{r} )  \alpha( \mathbf{q}, \mathbf{p}, \mathbf{r} )   \langle \hat{n}_{ \mathbf{q}-\mathbf{r}, \uparrow } \rangle :\hat{c}_{ \mathbf{p}-\mathbf{r}, \downarrow }^\dagger \hat{c}_{ \mathbf{p}-\mathbf{r}, \downarrow }^{\phantom{\dagger}}: \\
	&+1 \sum_{ \mathbf{q} \mathbf{p} \mathbf{r} } U( \mathbf{p}, \mathbf{r}, 0 )  \alpha( \mathbf{p}, \mathbf{q}, 0 )   \langle \hat{n}_{ \mathbf{q}, \downarrow } \rangle \langle \hat{n}_{ \mathbf{p}, \uparrow } \rangle \langle \hat{n}_{ \mathbf{p}, \uparrow } \rangle :\hat{c}_{ \mathbf{r}, \downarrow }^\dagger \hat{c}_{ \mathbf{r}, \downarrow }^{\phantom{\dagger}}: \\
	&-1 \sum_{ \mathbf{q} \mathbf{p} \mathbf{r} } U( \mathbf{q}, \mathbf{p}, \mathbf{r} )  \alpha( \mathbf{q}-\mathbf{r}, \mathbf{p}+\mathbf{r}, \mathbf{r} )   \langle \hat{n}_{ \mathbf{p}+\mathbf{r}, \downarrow } \rangle :\hat{c}_{ \mathbf{q}+\mathbf{r}, \uparrow }^\dagger \hat{c}_{ \mathbf{q}+\mathbf{r}, \uparrow }^{\phantom{\dagger}}: \\
	&-1 \sum_{ \mathbf{q} \mathbf{p} \mathbf{r} } U( \mathbf{q}, \mathbf{p}, \mathbf{r} )  \alpha( \mathbf{q}-\mathbf{r}, \mathbf{p}+\mathbf{r}, \mathbf{r} )   \langle \hat{n}_{ \mathbf{q}-\mathbf{r}, \uparrow } \rangle :\hat{c}_{ \mathbf{p}-\mathbf{r}, \downarrow }^\dagger \hat{c}_{ \mathbf{p}-\mathbf{r}, \downarrow }^{\phantom{\dagger}}: \\
	&+1 \sum_{ \mathbf{q} \mathbf{p} \mathbf{r} } U( \mathbf{p}, \mathbf{r}, 0 )  \alpha( \mathbf{p}, \mathbf{q}, 0 )   \langle \hat{n}_{ \mathbf{r}, \downarrow } \rangle \langle \hat{n}_{ \mathbf{p}, \uparrow } \rangle \langle \hat{n}_{ \mathbf{p}, \uparrow } \rangle :\hat{c}_{ \mathbf{q}, \downarrow }^\dagger \hat{c}_{ \mathbf{q}, \downarrow }^{\phantom{\dagger}}: \\
	&-1 \sum_{ \mathbf{q} \mathbf{p} \mathbf{r} } U( \mathbf{q}, \mathbf{p}, \mathbf{r} )  \alpha( \mathbf{q}-\mathbf{r}, \mathbf{p}+\mathbf{r}, \mathbf{r} )   \langle \hat{n}_{ \mathbf{q}, \uparrow } \rangle \langle \hat{n}_{ \mathbf{q}-\mathbf{r}, \uparrow } \rangle \langle \hat{n}_{ \mathbf{p}, \downarrow } \rangle :\hat{c}_{ \mathbf{p}+\mathbf{r}, \downarrow }^\dagger \hat{c}_{ \mathbf{p}+\mathbf{r}, \downarrow }^{\phantom{\dagger}}: \\
	&+1 \sum_{ \mathbf{q} \mathbf{p} \mathbf{r} } U( \mathbf{q}-\mathbf{r}, \mathbf{p}+\mathbf{r}, \mathbf{r} )  \alpha( \mathbf{q}, \mathbf{p}, \mathbf{r} )   \langle \hat{n}_{ \mathbf{q}-\mathbf{r}, \uparrow } \rangle \langle \hat{n}_{ \mathbf{p}, \downarrow } \rangle \langle \hat{n}_{ \mathbf{p}+\mathbf{r}, \downarrow } \rangle :\hat{c}_{ \mathbf{q}, \uparrow }^\dagger \hat{c}_{ \mathbf{q}, \uparrow }^{\phantom{\dagger}}: \\
	&+1 \sum_{ \mathbf{q} \mathbf{p} \mathbf{r} \mathbf{s} } U( \mathbf{q}, \mathbf{p}, 0 )  \alpha( \mathbf{r}, \mathbf{s}, 0 )   \langle \hat{n}_{ \mathbf{s}, \downarrow } \rangle \langle \hat{n}_{ \mathbf{q}, \uparrow } \rangle \langle \hat{n}_{ \mathbf{p}, \downarrow } \rangle :\hat{c}_{ \mathbf{r}, \uparrow }^\dagger \hat{c}_{ \mathbf{r}, \uparrow }^{\phantom{\dagger}}: \\
	&-1 \sum_{ \mathbf{q} \mathbf{p} \mathbf{r} } U( \mathbf{q}, \mathbf{r}, 0 )  \alpha( \mathbf{p}, \mathbf{r}, 0 )   \langle \hat{n}_{ \mathbf{r}, \downarrow } \rangle \langle \hat{n}_{ \mathbf{r}, \downarrow } \rangle \langle \hat{n}_{ \mathbf{q}, \uparrow } \rangle :\hat{c}_{ \mathbf{p}, \uparrow }^\dagger \hat{c}_{ \mathbf{p}, \uparrow }^{\phantom{\dagger}}: \\
	&+1 \sum_{ \mathbf{q} \mathbf{p} \mathbf{r} \mathbf{s} } U( \mathbf{q}, \mathbf{p}, 0 )  \alpha( \mathbf{r}, \mathbf{s}, 0 )   \langle \hat{n}_{ \mathbf{s}, \downarrow } \rangle \langle \hat{n}_{ \mathbf{r}, \uparrow } \rangle \langle \hat{n}_{ \mathbf{p}, \downarrow } \rangle :\hat{c}_{ \mathbf{q}, \uparrow }^\dagger \hat{c}_{ \mathbf{q}, \uparrow }^{\phantom{\dagger}}: \\
	&-1 \sum_{ \mathbf{q} \mathbf{p} \mathbf{r} } U( \mathbf{q}, \mathbf{r}, 0 )  \alpha( \mathbf{p}, \mathbf{r}, 0 )   \langle \hat{n}_{ \mathbf{r}, \downarrow } \rangle \langle \hat{n}_{ \mathbf{r}, \downarrow } \rangle \langle \hat{n}_{ \mathbf{p}, \uparrow } \rangle :\hat{c}_{ \mathbf{q}, \uparrow }^\dagger \hat{c}_{ \mathbf{q}, \uparrow }^{\phantom{\dagger}}: \\
	&-2 \sum_{ \mathbf{q} \mathbf{p} \mathbf{r} } U( \mathbf{q}, \mathbf{r}, 0 )  \alpha( \mathbf{p}, \mathbf{r}, 0 )   \langle \hat{n}_{ \mathbf{r}, \downarrow } \rangle \langle \hat{n}_{ \mathbf{q}, \uparrow } \rangle \langle \hat{n}_{ \mathbf{p}, \uparrow } \rangle :\hat{c}_{ \mathbf{r}, \downarrow }^\dagger \hat{c}_{ \mathbf{r}, \downarrow }^{\phantom{\dagger}}: \\
	&+1 \sum_{ \mathbf{q} \mathbf{p} \mathbf{r} } U( \mathbf{p}, \mathbf{q}, 0 )  \alpha( \mathbf{p}, \mathbf{r}, 0 )   \langle \hat{n}_{ \mathbf{r}, \downarrow } \rangle \langle \hat{n}_{ \mathbf{q}, \downarrow } \rangle :\hat{c}_{ \mathbf{p}, \uparrow }^\dagger \hat{c}_{ \mathbf{p}, \uparrow }^{\phantom{\dagger}}: \\
	&-1 \sum_{ \mathbf{q} \mathbf{p} \mathbf{r} } U( \mathbf{q}-\mathbf{r}, \mathbf{p}+\mathbf{r}, \mathbf{r} )  \alpha( \mathbf{q}, \mathbf{p}, \mathbf{r} )   \langle \hat{n}_{ \mathbf{p}, \downarrow } \rangle \langle \hat{n}_{ \mathbf{p}+\mathbf{r}, \downarrow } \rangle :\hat{c}_{ \mathbf{q}+\mathbf{r}, \uparrow }^\dagger \hat{c}_{ \mathbf{q}+\mathbf{r}, \uparrow }^{\phantom{\dagger}}: \\
	&-1 \sum_{ \mathbf{q} \mathbf{p} \mathbf{r} } U( \mathbf{q}-\mathbf{r}, \mathbf{p}+\mathbf{r}, \mathbf{r} )  \alpha( \mathbf{q}, \mathbf{p}, \mathbf{r} )   \langle \hat{n}_{ \mathbf{q}-\mathbf{r}, \uparrow } \rangle \langle \hat{n}_{ \mathbf{p}+\mathbf{r}, \downarrow } \rangle :\hat{c}_{ \mathbf{p}, \downarrow }^\dagger \hat{c}_{ \mathbf{p}, \downarrow }^{\phantom{\dagger}}: \\
	&+1 \sum_{ \mathbf{q} \mathbf{p} \mathbf{r} } U( \mathbf{p}, \mathbf{q}, 0 )  \alpha( \mathbf{p}, \mathbf{r}, 0 )   \langle \hat{n}_{ \mathbf{q}, \downarrow } \rangle \langle \hat{n}_{ \mathbf{p}, \uparrow } \rangle :\hat{c}_{ \mathbf{r}, \downarrow }^\dagger \hat{c}_{ \mathbf{r}, \downarrow }^{\phantom{\dagger}}: \\
	&+1 \sum_{ \mathbf{q} \mathbf{p} \mathbf{r} } U( \mathbf{p}, \mathbf{q}, 0 )  \alpha( \mathbf{p}, \mathbf{r}, 0 )   \langle \hat{n}_{ \mathbf{r}, \downarrow } \rangle \langle \hat{n}_{ \mathbf{p}, \uparrow } \rangle :\hat{c}_{ \mathbf{q}, \downarrow }^\dagger \hat{c}_{ \mathbf{q}, \downarrow }^{\phantom{\dagger}}: \\
	&-1 \sum_{ \mathbf{q} \mathbf{p} \mathbf{r} } U( \mathbf{q}-\mathbf{r}, \mathbf{p}+\mathbf{r}, \mathbf{r} )  \alpha( \mathbf{q}, \mathbf{p}, \mathbf{r} )   \langle \hat{n}_{ \mathbf{q}-\mathbf{r}, \uparrow } \rangle \langle \hat{n}_{ \mathbf{p}, \downarrow } \rangle :\hat{c}_{ \mathbf{p}-\mathbf{r}, \downarrow }^\dagger \hat{c}_{ \mathbf{p}-\mathbf{r}, \downarrow }^{\phantom{\dagger}}: \\
	&+1 \sum_{ \mathbf{q} \mathbf{p} \mathbf{r} \mathbf{s} } U( \mathbf{q}, \mathbf{p}, 0 )  \alpha( \mathbf{r}, \mathbf{s}, 0 )   \langle \hat{n}_{ \mathbf{r}, \uparrow } \rangle \langle \hat{n}_{ \mathbf{q}, \uparrow } \rangle \langle \hat{n}_{ \mathbf{p}, \downarrow } \rangle :\hat{c}_{ \mathbf{s}, \downarrow }^\dagger \hat{c}_{ \mathbf{s}, \downarrow }^{\phantom{\dagger}}: \\
	&+1 \sum_{ \mathbf{q} \mathbf{p} \mathbf{r} \mathbf{s} } U( \mathbf{q}, \mathbf{p}, 0 )  \alpha( \mathbf{r}, \mathbf{s}, 0 )   \langle \hat{n}_{ \mathbf{s}, \downarrow } \rangle \langle \hat{n}_{ \mathbf{r}, \uparrow } \rangle \langle \hat{n}_{ \mathbf{q}, \uparrow } \rangle :\hat{c}_{ \mathbf{p}, \downarrow }^\dagger \hat{c}_{ \mathbf{p}, \downarrow }^{\phantom{\dagger}}: \\
	&+1 \sum_{ \mathbf{q} \mathbf{p} \mathbf{r} } U( \mathbf{q}-\mathbf{r}, \mathbf{p}+\mathbf{r}, \mathbf{r} )  \alpha( \mathbf{q}, \mathbf{p}, \mathbf{r} )   \langle \hat{n}_{ \mathbf{q}, \uparrow } \rangle \langle \hat{n}_{ \mathbf{p}, \downarrow } \rangle \langle \hat{n}_{ \mathbf{p}+\mathbf{r}, \downarrow } \rangle :\hat{c}_{ \mathbf{q}-\mathbf{r}, \uparrow }^\dagger \hat{c}_{ \mathbf{q}-\mathbf{r}, \uparrow }^{\phantom{\dagger}}: \\
	&+1 \sum_{ \mathbf{q} \mathbf{p} \mathbf{r} } U( \mathbf{q}-\mathbf{r}, \mathbf{p}+\mathbf{r}, \mathbf{r} )  \alpha( \mathbf{q}, \mathbf{p}, \mathbf{r} )   \langle \hat{n}_{ \mathbf{q}, \uparrow } \rangle \langle \hat{n}_{ \mathbf{q}-\mathbf{r}, \uparrow } \rangle \langle \hat{n}_{ \mathbf{p}+\mathbf{r}, \downarrow } \rangle :\hat{c}_{ \mathbf{p}, \downarrow }^\dagger \hat{c}_{ \mathbf{p}, \downarrow }^{\phantom{\dagger}}: \\
	&-1 \sum_{ \mathbf{q} \mathbf{p} \mathbf{r} } U( \mathbf{p}, \mathbf{q}, 0 )  \alpha( \mathbf{p}, \mathbf{r}, 0 )   \langle \hat{n}_{ \mathbf{q}, \downarrow } \rangle \langle \hat{n}_{ \mathbf{p}, \uparrow } \rangle \langle \hat{n}_{ \mathbf{p}, \uparrow } \rangle :\hat{c}_{ \mathbf{r}, \downarrow }^\dagger \hat{c}_{ \mathbf{r}, \downarrow }^{\phantom{\dagger}}: \\
	&-1 \sum_{ \mathbf{q} \mathbf{p} \mathbf{r} } U( \mathbf{p}, \mathbf{q}, 0 )  \alpha( \mathbf{p}, \mathbf{r}, 0 )   \langle \hat{n}_{ \mathbf{r}, \downarrow } \rangle \langle \hat{n}_{ \mathbf{p}, \uparrow } \rangle \langle \hat{n}_{ \mathbf{p}, \uparrow } \rangle :\hat{c}_{ \mathbf{q}, \downarrow }^\dagger \hat{c}_{ \mathbf{q}, \downarrow }^{\phantom{\dagger}}: \\
	&-1 \sum_{ \mathbf{q} \mathbf{p} \mathbf{r} } U( \mathbf{q}-\mathbf{r}, \mathbf{p}+\mathbf{r}, \mathbf{r} )  \alpha( \mathbf{q}, \mathbf{p}, \mathbf{r} )   \langle \hat{n}_{ \mathbf{q}-\mathbf{r}, \uparrow } \rangle \langle \hat{n}_{ \mathbf{p}+\mathbf{r}, \downarrow } \rangle :\hat{c}_{ \mathbf{q}, \uparrow }^\dagger \hat{c}_{ \mathbf{q}, \uparrow }^{\phantom{\dagger}}: \\
	&+1 \sum_{ \mathbf{q} \mathbf{p} \mathbf{r} } U( \mathbf{q}, \mathbf{r}, 0 )  \alpha( \mathbf{p}, \mathbf{r}, 0 )   \langle \hat{n}_{ \mathbf{r}, \downarrow } \rangle \langle \hat{n}_{ \mathbf{q}, \uparrow } \rangle :\hat{c}_{ \mathbf{p}, \uparrow }^\dagger \hat{c}_{ \mathbf{p}, \uparrow }^{\phantom{\dagger}}: \\
	&+1 \sum_{ \mathbf{q} \mathbf{p} \mathbf{r} } U( \mathbf{q}, \mathbf{r}, 0 )  \alpha( \mathbf{p}, \mathbf{r}, 0 )   \langle \hat{n}_{ \mathbf{r}, \downarrow } \rangle \langle \hat{n}_{ \mathbf{p}, \uparrow } \rangle :\hat{c}_{ \mathbf{q}, \uparrow }^\dagger \hat{c}_{ \mathbf{q}, \uparrow }^{\phantom{\dagger}}: \\
	&+1 \sum_{ \mathbf{q} \mathbf{p} \mathbf{r} } U( \mathbf{q}, \mathbf{r}, 0 )  \alpha( \mathbf{p}, \mathbf{r}, 0 )   \langle \hat{n}_{ \mathbf{q}, \uparrow } \rangle \langle \hat{n}_{ \mathbf{p}, \uparrow } \rangle :\hat{c}_{ \mathbf{r}, \downarrow }^\dagger \hat{c}_{ \mathbf{r}, \downarrow }^{\phantom{\dagger}}: \\
	&-1 \sum_{ \mathbf{q} \mathbf{p} \mathbf{r} } U( \mathbf{q}-\mathbf{r}, \mathbf{p}+\mathbf{r}, \mathbf{r} )  \alpha( \mathbf{q}, \mathbf{p}, \mathbf{r} )   \langle \hat{n}_{ \mathbf{q}, \uparrow } \rangle \langle \hat{n}_{ \mathbf{p}+\mathbf{r}, \downarrow } \rangle :\hat{c}_{ \mathbf{q}+\mathbf{r}, \uparrow }^\dagger \hat{c}_{ \mathbf{q}+\mathbf{r}, \uparrow }^{\phantom{\dagger}}: \\
	&-1 \sum_{ \mathbf{q} \mathbf{p} \mathbf{r} } U( \mathbf{q}-\mathbf{r}, \mathbf{p}+\mathbf{r}, \mathbf{r} )  \alpha( \mathbf{q}, \mathbf{p}, \mathbf{r} )   \langle \hat{n}_{ \mathbf{q}, \uparrow } \rangle \langle \hat{n}_{ \mathbf{q}-\mathbf{r}, \uparrow } \rangle :\hat{c}_{ \mathbf{p}-\mathbf{r}, \downarrow }^\dagger \hat{c}_{ \mathbf{p}-\mathbf{r}, \downarrow }^{\phantom{\dagger}}: \\
	&+1 \sum_{ \mathbf{q} \mathbf{p} \mathbf{r} } U( \mathbf{q}-\mathbf{r}, \mathbf{p}+\mathbf{r}, \mathbf{r} )  \alpha( \mathbf{q}, \mathbf{p}, \mathbf{r} )   \langle \hat{n}_{ \mathbf{q}, \uparrow } \rangle \langle \hat{n}_{ \mathbf{q}-\mathbf{r}, \uparrow } \rangle \langle \hat{n}_{ \mathbf{p}, \downarrow } \rangle :\hat{c}_{ \mathbf{p}+\mathbf{r}, \downarrow }^\dagger \hat{c}_{ \mathbf{p}+\mathbf{r}, \downarrow }^{\phantom{\dagger}}: \\
	&-2 \sum_{ \mathbf{q} \mathbf{p} \mathbf{r} } U( \mathbf{p}, \mathbf{q}, 0 )  \alpha( \mathbf{p}, \mathbf{r}, 0 )   \langle \hat{n}_{ \mathbf{r}, \downarrow } \rangle \langle \hat{n}_{ \mathbf{q}, \downarrow } \rangle \langle \hat{n}_{ \mathbf{p}, \uparrow } \rangle :\hat{c}_{ \mathbf{p}, \uparrow }^\dagger \hat{c}_{ \mathbf{p}, \uparrow }^{\phantom{\dagger}}: \\
	&+2 \sum_{ \mathbf{q} \mathbf{p} \mathbf{r} } U( \mathbf{p}, \mathbf{r}, 0 )  \alpha( \mathbf{p}, \mathbf{q}, 0 )   \langle \hat{n}_{ \mathbf{r}, \downarrow } \rangle \langle \hat{n}_{ \mathbf{q}, \downarrow } \rangle \langle \hat{n}_{ \mathbf{p}, \uparrow } \rangle :\hat{c}_{ \mathbf{p}, \uparrow }^\dagger \hat{c}_{ \mathbf{p}, \uparrow }^{\phantom{\dagger}}: \\
	&-1 \sum_{ \mathbf{q} \mathbf{p} \mathbf{r} } U( \mathbf{q}, \mathbf{p}, \mathbf{r} )  \alpha( \mathbf{q}-\mathbf{r}, \mathbf{p}+\mathbf{r}, \mathbf{r} )   \langle \hat{n}_{ \mathbf{q}-\mathbf{r}, \uparrow } \rangle \langle \hat{n}_{ \mathbf{p}, \downarrow } \rangle \langle \hat{n}_{ \mathbf{p}+\mathbf{r}, \downarrow } \rangle :\hat{c}_{ \mathbf{q}, \uparrow }^\dagger \hat{c}_{ \mathbf{q}, \uparrow }^{\phantom{\dagger}}: \\
	&-1 \sum_{ \mathbf{q} \mathbf{p} \mathbf{r} \mathbf{s} } U( \mathbf{r}, \mathbf{s}, 0 )  \alpha( \mathbf{q}, \mathbf{p}, 0 )   \langle \hat{n}_{ \mathbf{s}, \downarrow } \rangle \langle \hat{n}_{ \mathbf{q}, \uparrow } \rangle \langle \hat{n}_{ \mathbf{p}, \downarrow } \rangle :\hat{c}_{ \mathbf{r}, \uparrow }^\dagger \hat{c}_{ \mathbf{r}, \uparrow }^{\phantom{\dagger}}: \\
	&+1 \sum_{ \mathbf{q} \mathbf{p} \mathbf{r} } U( \mathbf{p}, \mathbf{r}, 0 )  \alpha( \mathbf{q}, \mathbf{r}, 0 )   \langle \hat{n}_{ \mathbf{r}, \downarrow } \rangle \langle \hat{n}_{ \mathbf{r}, \downarrow } \rangle \langle \hat{n}_{ \mathbf{q}, \uparrow } \rangle :\hat{c}_{ \mathbf{p}, \uparrow }^\dagger \hat{c}_{ \mathbf{p}, \uparrow }^{\phantom{\dagger}}: \\
	&-1 \sum_{ \mathbf{q} \mathbf{p} \mathbf{r} } U( \mathbf{p}, \mathbf{r}, 0 )  \alpha( \mathbf{p}, \mathbf{q}, 0 )   \langle \hat{n}_{ \mathbf{r}, \downarrow } \rangle \langle \hat{n}_{ \mathbf{q}, \downarrow } \rangle :\hat{c}_{ \mathbf{p}, \uparrow }^\dagger \hat{c}_{ \mathbf{p}, \uparrow }^{\phantom{\dagger}}: \\
	&+1 \sum_{ \mathbf{q} \mathbf{p} \mathbf{r} } U( \mathbf{q}, \mathbf{p}, \mathbf{r} )  \alpha( \mathbf{q}-\mathbf{r}, \mathbf{p}+\mathbf{r}, \mathbf{r} )   \langle \hat{n}_{ \mathbf{p}, \downarrow } \rangle \langle \hat{n}_{ \mathbf{p}+\mathbf{r}, \downarrow } \rangle :\hat{c}_{ \mathbf{q}+\mathbf{r}, \uparrow }^\dagger \hat{c}_{ \mathbf{q}+\mathbf{r}, \uparrow }^{\phantom{\dagger}}: \\
	&+1 \sum_{ \mathbf{q} \mathbf{p} \mathbf{r} } U( \mathbf{q}, \mathbf{p}, \mathbf{r} )  \alpha( \mathbf{q}-\mathbf{r}, \mathbf{p}+\mathbf{r}, \mathbf{r} )   \langle \hat{n}_{ \mathbf{q}-\mathbf{r}, \uparrow } \rangle \langle \hat{n}_{ \mathbf{p}+\mathbf{r}, \downarrow } \rangle :\hat{c}_{ \mathbf{p}, \downarrow }^\dagger \hat{c}_{ \mathbf{p}, \downarrow }^{\phantom{\dagger}}: \\
	&-1 \sum_{ \mathbf{q} \mathbf{p} \mathbf{r} } U( \mathbf{p}, \mathbf{r}, 0 )  \alpha( \mathbf{p}, \mathbf{q}, 0 )   \langle \hat{n}_{ \mathbf{q}, \downarrow } \rangle \langle \hat{n}_{ \mathbf{p}, \uparrow } \rangle :\hat{c}_{ \mathbf{r}, \downarrow }^\dagger \hat{c}_{ \mathbf{r}, \downarrow }^{\phantom{\dagger}}: \\
	&-1 \sum_{ \mathbf{q} \mathbf{p} \mathbf{r} } U( \mathbf{p}, \mathbf{r}, 0 )  \alpha( \mathbf{p}, \mathbf{q}, 0 )   \langle \hat{n}_{ \mathbf{r}, \downarrow } \rangle \langle \hat{n}_{ \mathbf{p}, \uparrow } \rangle :\hat{c}_{ \mathbf{q}, \downarrow }^\dagger \hat{c}_{ \mathbf{q}, \downarrow }^{\phantom{\dagger}}: \\
	&+1 \sum_{ \mathbf{q} \mathbf{p} \mathbf{r} } U( \mathbf{q}, \mathbf{p}, \mathbf{r} )  \alpha( \mathbf{q}-\mathbf{r}, \mathbf{p}+\mathbf{r}, \mathbf{r} )   \langle \hat{n}_{ \mathbf{q}-\mathbf{r}, \uparrow } \rangle \langle \hat{n}_{ \mathbf{p}, \downarrow } \rangle :\hat{c}_{ \mathbf{p}-\mathbf{r}, \downarrow }^\dagger \hat{c}_{ \mathbf{p}-\mathbf{r}, \downarrow }^{\phantom{\dagger}}: \\
	&-1 \sum_{ \mathbf{q} \mathbf{p} \mathbf{r} \mathbf{s} } U( \mathbf{r}, \mathbf{s}, 0 )  \alpha( \mathbf{q}, \mathbf{p}, 0 )   \langle \hat{n}_{ \mathbf{s}, \downarrow } \rangle \langle \hat{n}_{ \mathbf{r}, \uparrow } \rangle \langle \hat{n}_{ \mathbf{p}, \downarrow } \rangle :\hat{c}_{ \mathbf{q}, \uparrow }^\dagger \hat{c}_{ \mathbf{q}, \uparrow }^{\phantom{\dagger}}: \\
	&+1 \sum_{ \mathbf{q} \mathbf{p} \mathbf{r} } U( \mathbf{p}, \mathbf{r}, 0 )  \alpha( \mathbf{q}, \mathbf{r}, 0 )   \langle \hat{n}_{ \mathbf{r}, \downarrow } \rangle \langle \hat{n}_{ \mathbf{r}, \downarrow } \rangle \langle \hat{n}_{ \mathbf{p}, \uparrow } \rangle :\hat{c}_{ \mathbf{q}, \uparrow }^\dagger \hat{c}_{ \mathbf{q}, \uparrow }^{\phantom{\dagger}}: \\
	&-1 \sum_{ \mathbf{q} \mathbf{p} \mathbf{r} \mathbf{s} } U( \mathbf{r}, \mathbf{s}, 0 )  \alpha( \mathbf{q}, \mathbf{p}, 0 )   \langle \hat{n}_{ \mathbf{s}, \downarrow } \rangle \langle \hat{n}_{ \mathbf{r}, \uparrow } \rangle \langle \hat{n}_{ \mathbf{q}, \uparrow } \rangle \langle \hat{n}_{ \mathbf{p}, \downarrow } \rangle : \hat{1} : \\
	&+1 \sum_{ \mathbf{q} \mathbf{p} \mathbf{r} } U( \mathbf{p}, \mathbf{r}, 0 )  \alpha( \mathbf{q}, \mathbf{r}, 0 )   \langle \hat{n}_{ \mathbf{r}, \downarrow } \rangle \langle \hat{n}_{ \mathbf{r}, \downarrow } \rangle \langle \hat{n}_{ \mathbf{q}, \uparrow } \rangle \langle \hat{n}_{ \mathbf{p}, \uparrow } \rangle : \hat{1} : \\
	&+1 \sum_{ \mathbf{q} \mathbf{p} \mathbf{r} } U( \mathbf{q}, \mathbf{r}, 0 )  \alpha( \mathbf{p}, \mathbf{r}, 0 )   \langle \hat{n}_{ \mathbf{r}, \downarrow } \rangle \langle \hat{n}_{ \mathbf{q}, \uparrow } \rangle \langle \hat{n}_{ \mathbf{p}, \uparrow } \rangle : \hat{1} : \\
	&-1 \sum_{ \mathbf{q} \mathbf{p} \mathbf{r} } U( \mathbf{q}-\mathbf{r}, \mathbf{p}+\mathbf{r}, \mathbf{r} )  \alpha( \mathbf{q}, \mathbf{p}, \mathbf{r} )   \langle \hat{n}_{ \mathbf{q}, \uparrow } \rangle \langle \hat{n}_{ \mathbf{q}-\mathbf{r}, \uparrow } \rangle \langle \hat{n}_{ \mathbf{p}+\mathbf{r}, \downarrow } \rangle : \hat{1} : \\
	&-1 \sum_{ \mathbf{q} \mathbf{p} \mathbf{r} } U( \mathbf{p}, \mathbf{r}, 0 )  \alpha( \mathbf{q}, \mathbf{r}, 0 )   \langle \hat{n}_{ \mathbf{r}, \downarrow } \rangle \langle \hat{n}_{ \mathbf{q}, \uparrow } \rangle \langle \hat{n}_{ \mathbf{p}, \uparrow } \rangle : \hat{1} : \\
	&+1 \sum_{ \mathbf{q} \mathbf{p} \mathbf{r} } U( \mathbf{q}, \mathbf{p}, \mathbf{r} )  \alpha( \mathbf{q}-\mathbf{r}, \mathbf{p}+\mathbf{r}, \mathbf{r} )   \langle \hat{n}_{ \mathbf{q}, \uparrow } \rangle \langle \hat{n}_{ \mathbf{q}-\mathbf{r}, \uparrow } \rangle \langle \hat{n}_{ \mathbf{p}+\mathbf{r}, \downarrow } \rangle : \hat{1} : \\
	&-1 \sum_{ \mathbf{q} \mathbf{p} \mathbf{r} } U( \mathbf{p}, \mathbf{r}, 0 )  \alpha( \mathbf{p}, \mathbf{q}, 0 )   \langle \hat{n}_{ \mathbf{r}, \downarrow } \rangle \langle \hat{n}_{ \mathbf{q}, \downarrow } \rangle \langle \hat{n}_{ \mathbf{p}, \uparrow } \rangle : \hat{1} : \\
	&+1 \sum_{ \mathbf{q} \mathbf{p} \mathbf{r} } U( \mathbf{q}, \mathbf{p}, \mathbf{r} )  \alpha( \mathbf{q}-\mathbf{r}, \mathbf{p}+\mathbf{r}, \mathbf{r} )   \langle \hat{n}_{ \mathbf{q}-\mathbf{r}, \uparrow } \rangle \langle \hat{n}_{ \mathbf{p}, \downarrow } \rangle \langle \hat{n}_{ \mathbf{p}+\mathbf{r}, \downarrow } \rangle : \hat{1} : \\
	&+1 \sum_{ \mathbf{q} \mathbf{p} \mathbf{r} } U( \mathbf{p}, \mathbf{r}, 0 )  \alpha( \mathbf{p}, \mathbf{q}, 0 )   \langle \hat{n}_{ \mathbf{r}, \downarrow } \rangle \langle \hat{n}_{ \mathbf{q}, \downarrow } \rangle \langle \hat{n}_{ \mathbf{p}, \uparrow } \rangle \langle \hat{n}_{ \mathbf{p}, \uparrow } \rangle : \hat{1} : \\
	&+1 \sum_{ \mathbf{q} \mathbf{p} \mathbf{r} } U( \mathbf{q}-\mathbf{r}, \mathbf{p}+\mathbf{r}, \mathbf{r} )  \alpha( \mathbf{q}, \mathbf{p}, \mathbf{r} )   \langle \hat{n}_{ \mathbf{q}-\mathbf{r}, \uparrow } \rangle \langle \hat{n}_{ \mathbf{p}+\mathbf{r}, \downarrow } \rangle : \hat{1} : \\
	&-1 \sum_{ \mathbf{q} \mathbf{p} \mathbf{r} } U( \mathbf{q}, \mathbf{p}, \mathbf{r} )  \alpha( \mathbf{q}-\mathbf{r}, \mathbf{p}+\mathbf{r}, \mathbf{r} )   \langle \hat{n}_{ \mathbf{q}, \uparrow } \rangle \langle \hat{n}_{ \mathbf{q}-\mathbf{r}, \uparrow } \rangle \langle \hat{n}_{ \mathbf{p}, \downarrow } \rangle \langle \hat{n}_{ \mathbf{p}+\mathbf{r}, \downarrow } \rangle : \hat{1} : \\
	&+1 \sum_{ \mathbf{q} \mathbf{p} \mathbf{r} } U( \mathbf{q}-\mathbf{r}, \mathbf{p}+\mathbf{r}, \mathbf{r} )  \alpha( \mathbf{q}, \mathbf{p}, \mathbf{r} )   \langle \hat{n}_{ \mathbf{q}, \uparrow } \rangle \langle \hat{n}_{ \mathbf{q}-\mathbf{r}, \uparrow } \rangle \langle \hat{n}_{ \mathbf{p}, \downarrow } \rangle \langle \hat{n}_{ \mathbf{p}+\mathbf{r}, \downarrow } \rangle : \hat{1} : \\
	&+1 \sum_{ \mathbf{q} \mathbf{p} \mathbf{r} \mathbf{s} } U( \mathbf{q}, \mathbf{p}, 0 )  \alpha( \mathbf{r}, \mathbf{s}, 0 )   \langle \hat{n}_{ \mathbf{s}, \downarrow } \rangle \langle \hat{n}_{ \mathbf{r}, \uparrow } \rangle \langle \hat{n}_{ \mathbf{q}, \uparrow } \rangle \langle \hat{n}_{ \mathbf{p}, \downarrow } \rangle : \hat{1} : \\
	&-1 \sum_{ \mathbf{q} \mathbf{p} \mathbf{r} } U( \mathbf{q}, \mathbf{p}, \mathbf{r} )  \alpha( \mathbf{q}-\mathbf{r}, \mathbf{p}+\mathbf{r}, \mathbf{r} )   \langle \hat{n}_{ \mathbf{q}-\mathbf{r}, \uparrow } \rangle \langle \hat{n}_{ \mathbf{p}+\mathbf{r}, \downarrow } \rangle : \hat{1} : \\
	&-1 \sum_{ \mathbf{q} \mathbf{p} \mathbf{r} } U( \mathbf{q}, \mathbf{r}, 0 )  \alpha( \mathbf{p}, \mathbf{r}, 0 )   \langle \hat{n}_{ \mathbf{r}, \downarrow } \rangle \langle \hat{n}_{ \mathbf{r}, \downarrow } \rangle \langle \hat{n}_{ \mathbf{q}, \uparrow } \rangle \langle \hat{n}_{ \mathbf{p}, \uparrow } \rangle : \hat{1} : \\
	&-1 \sum_{ \mathbf{q} \mathbf{p} \mathbf{r} } U( \mathbf{p}, \mathbf{q}, 0 )  \alpha( \mathbf{p}, \mathbf{r}, 0 )   \langle \hat{n}_{ \mathbf{r}, \downarrow } \rangle \langle \hat{n}_{ \mathbf{q}, \downarrow } \rangle \langle \hat{n}_{ \mathbf{p}, \uparrow } \rangle \langle \hat{n}_{ \mathbf{p}, \uparrow } \rangle : \hat{1} : \\
	&+1 \sum_{ \mathbf{q} \mathbf{p} \mathbf{r} } U( \mathbf{p}, \mathbf{q}, 0 )  \alpha( \mathbf{p}, \mathbf{r}, 0 )   \langle \hat{n}_{ \mathbf{r}, \downarrow } \rangle \langle \hat{n}_{ \mathbf{q}, \downarrow } \rangle \langle \hat{n}_{ \mathbf{p}, \uparrow } \rangle : \hat{1} : \\
	&-1 \sum_{ \mathbf{q} \mathbf{p} \mathbf{r} } U( \mathbf{q}-\mathbf{r}, \mathbf{p}+\mathbf{r}, \mathbf{r} )  \alpha( \mathbf{q}, \mathbf{p}, \mathbf{r} )   \langle \hat{n}_{ \mathbf{q}-\mathbf{r}, \uparrow } \rangle \langle \hat{n}_{ \mathbf{p}, \downarrow } \rangle \langle \hat{n}_{ \mathbf{p}+\mathbf{r}, \downarrow } \rangle : \hat{1} :
\end{align*}


\section{After the flow}

We track the the residual offdiagonality (ROD)
\begin{equation}
  \mathrm{ROD} (\ell) \coloneqq \sqrt{\sum_{\mathbf{p} \mathbf{q} \mathbf{r}} \left( U_{\uparrow \downarrow}^2 (\mathbf{p}, \mathbf{q}, \mathbf{r}; \ell) + U_{\parallel}^2 (\mathbf{p}, \mathbf{q}, \mathbf{r}; \ell) \right)}
\end{equation}
across the flow.
At some finite $\ell_0$ it will be minimized, though not zero.
At this point, we stop the flow and use the effective Hamiltonian $H_\mathrm{eff} \coloneqq H(\ell_0)$ for the remainder.
Besides the global constant, which we do not care about, it reads
\begin{align}
  H_\mathrm{eff} &= \sum_{\mathbf{p} \sigma} \tilde{\varepsilon}(\mathbf{p}; \ell_0) \hat{c}_{ \mathbf{p}, \sigma }^\dagger \hat{c}_{ \mathbf{p}, \sigma }^{\phantom{\dagger}} \\
    &+ \sum_{ \mathbf{p} \mathbf{q} \mathbf{r} } \sum_{ \sigma } U_{ \uparrow \downarrow }( \mathbf{p}, \mathbf{q}, \mathbf{r}; \ell_0 )   \hat{c}_{ \mathbf{p}, \sigma }^\dagger  \hat{c}_{ \mathbf{q}, -\sigma }^\dagger  \hat{c}_{ \mathbf{q}-\mathbf{r}, -\sigma }^{\phantom{\dagger}} \hat{c}_{ \mathbf{p}+\mathbf{r}, \sigma }^{\phantom{\dagger}} \\
    &+ \sum_{ \mathbf{p} \mathbf{q} \mathbf{r} } \sum_{ \sigma } U_{ \parallel }( \mathbf{p}, \mathbf{q}, \mathbf{r}; \ell_0 )   \hat{c}_{ \mathbf{p}, \sigma }^\dagger  \hat{c}_{ \mathbf{q}, \sigma }^\dagger  \hat{c}_{ \mathbf{q}-\mathbf{r}, \sigma }^{\phantom{\dagger}} \hat{c}_{ \mathbf{p}+\mathbf{r}, \sigma }^{\phantom{\dagger}} .
\end{align}
The goal is to apply a mean-field approximation to this effective Hamiltonian.
Besides the number operators, we will retain the expectation values relevant to superconductivity (SC) and density waves (DW), i.e,
\begin{align}
  \langle n_{\mathbf{p},\sigma} \rangle &\coloneqq \langle \hat{c}_{ \mathbf{p}, \sigma }^\dagger \hat{c}_{ \mathbf{p}, \sigma }^{\phantom{\dagger}} \rangle \\
  \langle g_{\mathbf{p},\sigma} \rangle &\coloneqq \langle \hat{c}_{ \mathbf{p}, \sigma }^\dagger \hat{c}_{ \mathbf{p} + \mathbf{Q}, \sigma }^{\phantom{\dagger}} \rangle \\
  \langle f_\mathbf{p} \rangle &\coloneqq \langle \hat{c}_{ -\mathbf{p}, \downarrow }^{\phantom{\dagger}} \hat{c}_{ \mathbf{p}, \uparrow }^{\phantom{\dagger}} \rangle.
\end{align}
The mean-field Hamiltonian thus reads
\begin{equation}
  H_\mathrm{MF} = \sum_{\mathbf{p}}' \mathbf{\Psi}_{\mathbf{p}}^\dagger
  \begin{pmatrix}
    \varepsilon^\uparrow (\mathbf{p}) & \bar{\Delta}_\mathrm{DW}^{\uparrow} (\mathbf{p}) & \bar{\Delta}_\mathrm{SC} (\mathbf{p}) & 0 \\
    \Delta_\mathrm{DW}^{\uparrow} (\mathbf{p}) & \varepsilon^\uparrow (\mathbf{p} + \mathbf{Q}) & 0 & \bar{\Delta}_\mathrm{SC} (\mathbf{p} + \mathbf{Q}) \\
    \Delta_\mathrm{SC} (\mathbf{p}) & 0 & - \varepsilon^\downarrow (- \mathbf{p}) & -\bar{\Delta}_\mathrm{DW}^{\downarrow} (-\mathbf{p}) \\
    0 & \Delta_\mathrm{SC} (\mathbf{p} + \mathbf{Q}) & -\Delta_\mathrm{DW}^{\downarrow} (-\mathbf{p}) & - \varepsilon^\downarrow (- \mathbf{p} - \mathbf{Q})
  \end{pmatrix}
  \mathbf{\Psi}_{\mathbf{p}}
\end{equation}
with the spinors
\begin{equation}
  \mathbf{\Psi}_{\mathbf{p}}^\dagger \coloneqq \begin{pmatrix}
    \hat{c}_{ \mathbf{p}, \uparrow }^\dagger & \hat{c}_{ \mathbf{p} + \mathbf{Q}, \uparrow }^\dagger & \hat{c}_{ -\mathbf{p}, \downarrow }^{\phantom{\dagger}} & \hat{c}_{ -\mathbf{p} - \mathbf{Q}, \downarrow }^{\phantom{\dagger}}
  \end{pmatrix}
  \quad\text{ and }\quad 
  \mathbf{\Psi}_{\mathbf{p}} \coloneqq \begin{pmatrix}
    \hat{c}_{ \mathbf{p}, \uparrow }^{\phantom{\dagger}} \\ \hat{c}_{ \mathbf{p} + \mathbf{Q}, \uparrow }^{\phantom{\dagger}} \\ \hat{c}_{ -\mathbf{p}, \downarrow }^{\dagger} \\ \hat{c}_{ -\mathbf{p} - \mathbf{Q}, \downarrow }^{\dagger}
  \end{pmatrix}.
\end{equation}
The primed sum indicates summation over the magnetic Brillouin zone to avoid double counting.
The mean-field parameters
\begin{align}
  \Delta_\mathrm{SC}(\mathbf{p}) &= 2 \sum_{\mathbf{q}} U_{\uparrow \downarrow} (\mathbf{p}, \mathbf{-p}, \mathbf{q} - \mathbf{p}; \ell_0) \langle f_{\mathbf{p}} \rangle \\
  \Delta_\mathrm{DW}^{\sigma} (\mathbf{p}) &= 2 \sum_{\mathbf{q}} U_{\uparrow \downarrow} (\mathbf{p}, \mathbf{q}, \mathbf{Q}; \ell_0) \langle g_{\mathbf{q}, -\sigma} \rangle \\
    &+ 2 \sum_{\mathbf{q}} \Big[  U_{\parallel} (\mathbf{p}, \mathbf{q}, \mathbf{Q}; \ell_0)  
     - U_{\parallel} (\mathbf{p}, \mathbf{q}, \mathbf{q} - \mathbf{p} + \mathbf{Q}; \ell_0) \Big] \langle g_{\mathbf{q}, \sigma} \rangle \\
  \varepsilon_\mathrm{I}^{\sigma}(\mathbf{p}) &= 2 \sum_{\mathbf{q}} U_{\uparrow \downarrow} (\mathbf{p}, \mathbf{q}, 0; \ell_0) \langle n_{\mathbf{q}, -\sigma} \rangle \\
    &+ 2 \sum_{\mathbf{q}} \Big[  U_{\parallel} (\mathbf{p}, \mathbf{q}, 0; \ell_0)  
     - U_{\parallel} (\mathbf{p}, \mathbf{q}, \mathbf{q} - \mathbf{p}; \ell_0) \Big] \langle n_{\mathbf{q}, \sigma} \rangle \\
  \varepsilon^\sigma (\mathbf{p}) &= \tilde{\varepsilon}(\mathbf{p}; \ell_0) + \varepsilon_\mathrm{I}^{\sigma}(\mathbf{p}).
\end{align}
must be computed self-consistently.
Note that $\mathbf{Q} = (\pi, \pi)^\mathrm{T}$ and that therefore $g_{\mathbf{q}, \sigma} = g_{\mathbf{q} + \mathbf{Q}, \sigma}^\dagger$.
We abused the following relations
\begin{align}
  U_{\alpha} (\mathbf{p}, \mathbf{q}, \mathbf{r}) &= U_{\alpha} (\mathbf{q}, \mathbf{p}, -\mathbf{r}) \\
  U_{\alpha} (\mathbf{p}, \mathbf{q}, \mathbf{Q}) &= U_{\alpha} (\mathbf{q}, \mathbf{p}, \mathbf{Q}).
\end{align}
The order parameter of the superconducting phase is $\Delta_\mathrm{SC}(\mathbf{p})$, the order parameter of the charge-ordered phase is 
\begin{equation}
  \Delta_\mathrm{CDW} (\mathbf{p}) \coloneqq \frac{ \Delta_\mathrm{DW}^{\uparrow}(\mathbf{p}) + \Delta_\mathrm{DW}^{\downarrow}(\mathbf{p}) }{2},
\end{equation}
and the one of the antiferromagnetic phase is
\begin{equation}
  \Delta_\mathrm{AFM} (\mathbf{p}) \coloneqq \frac{ \Delta_\mathrm{DW}^{\uparrow}(\mathbf{p}) - \Delta_\mathrm{DW}^{\downarrow}(\mathbf{p}) }{2}.
\end{equation}

The next step is to compute the pertinent expectation values.
To do so, let $\mathcal{U}_{\mathbf{p}}$ be the unitary transformation that diagonalizes $H_\mathrm{MF}$, i.e.,
\begin{align}
  H_\mathrm{MF} = &\sum_{\mathbf{p}} \mathbf{\Psi}_{\mathbf{p}}^\dagger \mathcal{U}_{\mathbf{p}} \ \mathrm{diag}(E_1, E_2, E_3, E_4)\ \mathcal{U}_{\mathbf{p}}^\dagger \mathbf{\Psi}_{\mathbf{p}} \\
    \coloneqq &\sum_{\mathbf{p}} \mathbf{\Phi}_{\mathbf{p}}^\dagger \ \mathrm{diag}(E_1, E_2, E_3, E_4)\ \mathbf{\Phi}_{\mathbf{p}}.
\end{align}
Then, we know that
\begin{equation}
    \braket{\mathbf{\Phi}_{\mathbf{p}} \mathbf{\Phi}_{\mathbf{p}}^\dagger}_{ij} = \delta_{ij} (1 - F(\beta E_i))\,,
\end{equation}
where $F(x)$ is the Fermi function and $\beta$ the inverse temperature.
Transforming back into the original basis yields the desired expectation values
\begin{equation}
  \rho_{\mathbf{p}} = \braket{\mathbf{\Psi}_{\mathbf{p}} \mathbf{\Psi}_{\mathbf{p}}^\dagger} = \mathcal{U}_{\mathbf{p}} \braket{\mathbf{\Phi}_{\mathbf{p}} \mathbf{\Phi}_{\mathbf{p}}^\dagger} \mathcal{U}_{\mathbf{p}}^\dagger 
  = \begin{pmatrix}
        1 - \braket{n_{\mathbf{p},\uparrow}} & -\braket{g_{\mathbf{p},\uparrow}^\dagger} & -\braket{f_{\mathbf{p}}} & 0 \\
        -\braket{g_{\mathbf{p},\uparrow}} & 1 - \braket{n_{\mathbf{p}+\mathbf{Q},\uparrow}} & 0 & -\braket{f_{\mathbf{p}+\mathbf{Q}}} \\
        -\braket{f_{\mathbf{p}}^\dagger} & 0 & \braket{n_{-k\downarrow}} & \braket{g_{-\mathbf{p},\downarrow}} \\
        0 & -\braket{f_{\mathbf{p}+\mathbf{Q}}^\dagger} & \braket{g_{-\mathbf{p},\downarrow}^\dagger} & \braket{n_{-\mathbf{p}-\mathbf{Q},\downarrow}}
    \end{pmatrix}.
\end{equation}

\end{document}